% ===========================================================================
%  纵列双发矢量推力飞行器：构型原理、动力学建模与控制特性分析
%  Tandem Twin Vector-Thrust Aircraft:
%  Configuration Principle, Dynamic Modeling, and Control Characteristics
% ===========================================================================
%  编译：xelatex 或 lualatex（UTF-8 中文支持）
% ===========================================================================

\documentclass[a4paper,12pt]{ctexart}

% ---------- 页面布局 ----------
\usepackage[margin=2.5cm]{geometry}
\usepackage{setspace}
\onehalfspacing

% ---------- 数学 ----------
\usepackage{amsmath,amssymb,amsthm}
\usepackage{mathtools}
\usepackage{bm}

\theoremstyle{plain}
\newtheorem{theorem}{定理}[section]
\newtheorem{lemma}[theorem]{引理}
\newtheorem{proposition}[theorem]{命题}
\theoremstyle{definition}
\newtheorem{definition}[theorem]{定义}
\newtheorem{remark}[theorem]{评注}

% ---------- 图表 ----------
\usepackage{graphicx}
\usepackage{float}
\usepackage{subcaption}
\usepackage{tikz}
\usepackage{pgfplots}
\pgfplotsset{compat=1.18}
\usetikzlibrary{arrows.meta,shapes,positioning,calc,patterns,angles,quotes,3d,backgrounds,fit}
\usepackage{booktabs}
\usepackage{longtable}
\usepackage{array}

% ---------- 超链接与引用 ----------
\usepackage[hidelinks]{hyperref}
\usepackage{cleveref}

% ---------- 代码 ----------
\usepackage{listings}
\lstset{
  basicstyle=\ttfamily\small,
  breaklines=true,
  numbers=left,
  numberstyle=\tiny,
  frame=single,
  columns=flexible,
}

% ---------- 其他 ----------
\usepackage{enumitem}
\usepackage{tabularx}

\title{\textbf{纵列双发矢量推力飞行器：\\构型原理、动力学建模与控制特性分析}}
\author{LShang}
\date{\today}

\begin{document}

\maketitle

\begin{center}
\rule{0.6\textwidth}{0.5pt}\\[6pt]
{\small
\textbf{摘要}：本文系统阐述一种新型固定翼飞行器构型——纵列双发正交单轴矢量推力布局的理论基础。
该构型以前后纵列双电机、正交单轴摆座、反向旋转螺旋桨及差速反扭滚转控制为核心特征，通过推力矢量化实现三轴控制力矩的解耦生成。
本文从第一性原理出发，完整推导推进系统的集总参数模型（基于桨盘动量理论、叶素理论与量纲分析的相似律）、
推力矢量到机体系的六自由度力/力矩映射（含摆座运动学、力臂生成与耦合项分离）、
整机六自由度刚体动力学方程（牛顿-欧拉方程与四元数姿态运动学）、
以及级联式增稳控制律的设计框架（包含角速率阻尼、姿态比例-积分反馈与角速度闭环模式）。
文中对控制效能矩阵的对角化性质、耦合项的量级评估、各飞行模态的物理机理和SAS参数整定原则进行了深入分析，
并系统比较了本构型与传统舵面控制及通用过驱动推力矢量平台的区别。
全文不绑定具体飞行器物理参数，重在揭示构型本身的动力学本质与控制特性。
}
\end{center}

\rule{\textwidth}{0.5pt}

\tableofcontents

\newpage

% ===========================================================================
\section{引言}
\label{sec:intro}
% ===========================================================================

\subsection{研究背景与问题定位}

固定翼飞行器的姿态控制传统上依赖气动舵面——副翼控制滚转、升降舵控制俯仰、方向舵控制偏航。
舵面偏转改变当地翼型的有效弯度，产生附加气动力，经力臂形成控制力矩。这一范式在常规飞行包线内成熟可靠，
经历了一个多世纪的工程验证，现代民航客机和高性能战斗机均以此为基础\cite{stevens2015aircraft}。

然而，气动舵面控制存在若干无法回避的物理局限。其一，舵面控制效能随动压（$\propto V^2$）衰减——
在起飞滑跑初期和着陆拉平阶段，舵面偏转产生的力矩可能不足以克服扰动或执行机动指令。
其二，舵面偏转不可避免地产生附加阻力，降低巡航效率。
其三，大迎角下翼面流动分离导致舵面浸没在分离区内，舵效骤降甚至出现反效——这是失速/尾旋事故的重要诱因。

多旋翼飞行器的兴起提供了另一种思路：通过电机转速差动直接产生控制力矩，彻底摆脱对气动舵面的依赖\cite{mahony2012multirotor}。
多旋翼在悬停和低速段的操控性远超固定翼——每个旋翼既是升力源又是控制执行器，控制带宽不受动压限制。
然而，多旋翼的续航和速度性能远逊于固定翼，这是由桨盘载荷与气动效率的根本差异决定的。
固定翼的升阻比通常在10--20量级，而旋翼悬停的等效升阻比（figure of merit）在0.7--0.8——
前者以向前飞行的动能换取升力，后者以向下加速空气的动量换取升力，能量效率的差距是物理性的。

矢量推力（thrust vectoring）技术作为第三种范式，通过改变发动机喷流（或螺旋桨滑流）方向来产生直接力矩。
该技术在军用航空领域已有成熟应用——如F-22的二元推力矢量喷管和Su-35的三元推力矢量喷管，
能够在失速后（post-stall）保持可控性，完成常规飞机无法实现的超机动动作\cite{stevens2015aircraft}。
其物理本质是：不依赖气动面与来流的相互作用，而是通过动量通量的方向偏转直接产生力和力矩。

将矢量推力概念引入小型电推进飞行器，结合多旋翼的差速控制思想（利用反扭矩差动控制滚转），
可望实现一种全新的构型范式——纵列双发矢量推力布局。该构型的核心创新在于：
\textbf{通过正交摆轴的几何设计，使每个控制通道由单一执行器主导，从构型层面实现通道间的天然解耦}，
从而规避了通用过驱动矢量推力平台所需的在线控制分配优化问题\cite{johansen2013control}。

\subsection{现有构型的分类与定位}

为理解本构型的创新点，有必要对现有矢量推力/垂直起降构型作简要分类。从推力方向的可控性角度，
可将电推进飞行器分为以下几类：

\textbf{固定推力方向}：所有常规固定翼和多旋翼均属此类。推力方向固连于机体，姿态控制通过
气动舵面（固定翼）或差速（多旋翼）实现。结构简单、可靠性高，但低速段控制权限受限（固定翼）
或巡航效率低（多旋翼）。

\textbf{倾转旋翼/倾转机翼}：V-22 Osprey和AW609为代表。旋翼短舱或整个机翼可倾转，
在垂直起降和高速巡航之间切换。倾转机构复杂、重量代价大，且倾转过渡段的控制律设计
（增益调度或非线性控制）具有相当的理论与工程难度\cite{kang2009control}。

\textbf{通用多自由度矢量推力}：每个推进器可在两个方向独立倾斜（如万向节悬挂），
$n$个推进器产生极大的执行空间冗余，控制分配为欠定问题\cite{seguigasco2014novel}。
这类构型灵活性极高，理论上可实现任意方向平动和姿态的独立控制，
但代价是机械复杂度高（每个推进器需两套伺服机构）、控制分配需要在线求解优化问题、
且多个推进器之间的气动干扰效应复杂\cite{zhu2025experimental}。

\textbf{纵列双发正交单轴矢量推力（本文构型）}：每个推进器仅有一个摆动自由度，
两摆轴正交，加上差速通道，总计四个执行器恰好控制四个自由度（推力+三轴力矩）。
摆轴正交的几何设计使控制效能矩阵退化为对角阵——从几何层面消除了通道间的一阶耦合。
该构型可视为上述三类之间的“甜点”：兼具矢量推力的全速域可控性、多旋翼的差速滚转机制、
以及固定翼的气动效率，同时以最少的执行器数量实现了完整的控制能力。

\subsection{构型概念}

本文研究的构型由以下核心要素构成：

\begin{enumerate}[label=(\arabic*)]
    \item \textbf{纵列双发}：两台电机沿机体纵轴（$x_b$）前后布置。前电机为拉力式
          （tractor），螺旋桨位于重心前方；尾电机为推进式（pusher），
          螺旋桨位于重心后方。双发沿纵轴布置最大限度地利用了机身长度作为摆动力臂；
    \item \textbf{反向旋转}：前后螺旋桨旋向相反——前电机正转（CW，从后向前看），
          尾电机反转（CCW），使额定工况下反扭矩在机体系$x$轴上相互抵消。
          反向旋转的另一好处是消除P-factor（非对称桨叶载荷）和陀螺进动的一阶效应；
    \item \textbf{正交单轴摆座}：前电机绕机体$z_b$轴摆动（偏航主控），尾电机绕
          机体$y_b$轴摆动（俯仰主控）。两摆轴正交，使前电机的摆动仅产生侧向力
          和偏航力矩（不产生$z$向力和俯仰力矩），尾电机的摆动仅产生$z$向力和俯仰力矩
          （不产生$y$向力和偏航力矩）——此即天然解耦的几何根源；
    \item \textbf{差速反扭滚转}：打破两发电机的反扭矩平衡（通过差速指令$\Delta\omega$
          使一台加速、另一台减速），产生净滚转力矩。这一机制本质上是共轴直升机
          扭矩差控制原理在纵列布局上的推广。
\end{enumerate}

该构型的核心设计哲学是：\textbf{用几何设计换取算法简洁性}。
在通用矢量推力平台中，控制分配是$6 \times 3n$的欠定优化问题，需要在线求解伪逆或二次规划；
在本构型中，正交摆轴的几何约束将执行空间维数恰好降至控制目标维数，
控制分配退化为简单的直接映射（三通道独立增益），无需任何在线优化计算。

\begin{figure}[H]
\centering
\begin{tikzpicture}[>=Stealth, scale=0.95,
  body/.style={draw=black!70, fill=black!8, rounded corners=2pt},
  motor/.style={draw=black, fill=black!20, minimum size=8mm, circle},
  arrow/.style={->, thick, >={Stealth[length=3mm]}},
  dim/.style={<->, thin, draw=black!50},
  annot/.style={font=\footnotesize\sffamily},
]
  % ====== 俯视图 (Top View) ======
  \begin{scope}[shift={(0,0)}]
    % 机体
    \draw[body] (-2.2,-1.0) rectangle (2.8,1.0);
    % 机翼
    \draw[body] (-0.8,-2.5) rectangle (0.6,2.5);
    % CG
    \fill (0,0) circle (1.8pt) node[annot, below right=1pt] {CG};
    % 前电机（拉力式）
    \node[motor] (fm) at (1.8,0) {};
    \draw[fill=black!10] (fm) circle (5mm);
    \draw[->] (fm) ++(35:3mm) arc (35:395:3mm);
    \node[annot] at (1.8,0.85) {前电机};
    \node[annot] at (1.8,0.55) {CW};
    % 尾电机（推进式）
    \node[motor] (rm) at (-1.6,0) {};
    \draw[fill=black!10] (rm) circle (5mm);
    \draw[->] (rm) ++(215:3mm) arc (215:-145:3mm);
    \node[annot] at (-1.6,0.85) {尾电机};
    \node[annot] at (-1.6,0.55) {CCW};
    % x_b 轴
    \draw[arrow] (-2.6,0) -- (3.2,0) node[annot, right] {$x_b$};
    % y_b 轴
    \draw[arrow] (0,-2.8) -- (0,2.8) node[annot, above] {$y_b$};
    % 力臂标注
    \draw[dim] (0,-1.5) -- node[annot, fill=white] {$a$} (1.8,-1.5);
    \draw[dim] (-1.6,-1.5) -- node[annot, fill=white] {$b$} (0,-1.5);
    % 前摆轴（z_b，垂直纸面）
    \node[annot, blue!60!black] at (1.8,-0.6) {$\delta_f$绕$z_b$};
    \draw[blue!60!black, ->] (1.8,-0.4) arc (-90:60:4mm);
    % 偏航力矩示意
    \draw[arrow, red!60!black] (0.3,1.2) arc (90:270:3mm and 5mm) node[annot, below=0pt, red!60!black] {$M_z$(偏航)};
  \end{scope}

  % ====== 侧视图 (Side View) ======
  \begin{scope}[shift={(7.5,0)}]
    % 机体
    \draw[body] (-2.2,-0.6) rectangle (2.8,0.8);
    % 机翼（侧视仅见剖面）
    \draw[body] (-0.2,-0.6) -- (-0.2,-1.3) .. controls (0.1,-1.3) and (0.3,-1.3) .. (0.4,-0.6);
    % CG
    \fill (0,0) circle (1.8pt) node[annot, below left=1pt] {CG};
    % 前电机
    \node[motor] (fms) at (1.8,0) {};
    \draw[fill=black!10] (fms) circle (4.5mm);
    \node[annot] at (1.8,0.8) {前};
    % 尾电机
    \node[motor] (rms) at (-1.6,0) {};
    \draw[fill=black!10] (rms) circle (4.5mm);
    \node[annot] at (-1.6,0.8) {尾};
    % 尾摆示意（δ_t > 0: 推力向下）
    \draw[arrow, red!60!black, dashed] (rms) -- ++(-25:1.3cm) node[annot, right=-1pt, red!60!black] {$\bm{T}_t$};
    \draw[arrow, red!60!black, dashed] (rms) -- ++(0:1.3cm);
    \draw[red!60!black] (rms) ++(0:7mm) arc (0:-25:7mm);
    \node[annot, red!60!black] at (-1.2,-0.5) {$\delta_t$};
    % 俯仰力矩标注
    \node[annot, red!60!black] at (-0.3,1.1) {$M_y$(俯仰)};
    \draw[->, red!60!black] (-0.3,0.9) arc (0:-270:3mm and 2mm);
    % x_b 轴
    \draw[arrow] (-2.6,0) -- (3.2,0) node[annot, right] {$x_b$};
    % z_b 轴 (NED: 向下为正)
    \draw[arrow] (0,1.6) -- (0,-2.0) node[annot, below] {$z_b$（向下）};
    % 力臂
    \draw[dim] (0,1.2) -- node[annot, fill=white] {$a$} (1.8,1.2);
    \draw[dim] (-1.6,1.2) -- node[annot, fill=white] {$b$} (0,1.2);
    % 尾摆轴（y_b = 垂直纸面）
    \node[annot, blue!60!black] at (-1.6,-1.1) {$\delta_t$绕$y_b$};
  \end{scope}

  \node[annot] at (3.75,-2.8) {(a) 俯视图（$xy$ 平面）};
  \node[annot] at (11.25,-2.8) {(b) 侧视图（$xz$ 平面）};
\end{tikzpicture}
\caption{纵列双发矢量推力飞行器构型布局示意。前电机（拉力式，CW）绕$z_b$轴摆动，控制偏航通道；
尾电机（推进式，CCW）绕$y_b$轴摆动，控制俯仰通道。零摆角时两电机推力线均通过质心CG。}
\label{fig:configuration}
\end{figure}

\subsection{本文结构}

本文按从底层物理到顶层控制的逻辑链展开：
第\ref{sec:propulsion}节建立推进系统的集总参数模型，从桨盘动量理论出发推导推力与反扭矩的相似律，
并分析电机动态和陀螺效应的物理机制；
第\ref{sec:mapping}节导出摆座运动学，给出完整的六维力/力矩映射，分离主控项与耦合项；
第\ref{sec:dynamics}节建立整机六自由度刚体动力学方程，包括牛顿-欧拉方程、四元数姿态运动学及数值积分策略；
第\ref{sec:allocation}节分析控制效能矩阵的对角化性质，评估耦合项的量级，并与通用推力矢量平台进行对比；
第\ref{sec:stability}节进行小扰动线性化，分析纵向和横航向各飞行模态的物理特性；
第\ref{sec:control}节设计级联式增稳控制律框架，给出各通道SAS律的整定原则；
第\ref{sec:discussion}节系统讨论构型的优势、已知局限、设计权衡及扩展方向。

% ===========================================================================
\section{推进系统建模}
\label{sec:propulsion}
% ===========================================================================

推进系统的建模是整机动力学分析的起点。本节从最基本的流体力学原理出发，
建立推力与反扭矩的数学描述，并分析电机动态响应和旋转部件的陀螺效应。

\subsection{桨盘动量理论}

动量理论（momentum theory）或称为Froude-Rankine理论，是分析螺旋桨/旋翼性能的最基本工具。
该理论将螺旋桨理想化为无限薄桨盘（actuator disk），忽略桨叶的具体几何形状，
将问题简化为控制体上的质量、动量和能量守恒\cite{leishman2006helicopter}。

\begin{figure}[H]
\centering
\begin{tikzpicture}[>=Stealth, scale=1.1,
  stream/.style={draw=black!40, thick},
  arrow/.style={->, thick, >={Stealth[length=2.5mm]}},
  annot/.style={font=\footnotesize\sffamily},
]
  % 流管上边界
  \draw[stream] (0,0) .. controls (2.5,1.2) and (3.5,1.2) .. (5.5,0.6);
  \draw[stream] (0,0) .. controls (2.5,-1.2) and (3.5,-1.2) .. (5.5,-0.6);
  % 截面0（远前方）
  \draw[stream, fill=black!5] (-0.1,0) -- (-0.1,-0.6) .. controls (0.5,-0.8) and (0.5,0.8) .. (-0.1,0.6) -- cycle;
  \draw[<->] (-0.1,0.7) -- node[annot, fill=white] {$A_0$} (-0.1,-0.7);
  % 桨盘
  \draw[stream, fill=black!12, line width=1.4pt] (2.5,0.65) -- (2.5,-0.65) node[annot, below=4pt] {桨盘};
  \draw[<->] (2.5,0.75) -- node[annot, fill=white] {$A$} (2.5,-0.75);
  % 截面2（远后方）
  \draw[stream, fill=black!5] (5.6,0) -- (5.6,-0.35) .. controls (5.0,-0.5) and (5.0,0.5) .. (5.6,0.35) -- cycle;
  % 速度标注
  \draw[arrow, blue!60!black] (-0.8,0) -- node[annot, above] {$v_0$} (0.3,0);
  \draw[arrow, blue!60!black] (1.5,0) -- node[annot, above] {$v_p=v_0+v_i$} (3.2,0);
  \draw[arrow, blue!60!black] (3.8,0) -- node[annot, above] {$v_2=v_0+2v_i$} (5.8,0);
  % 压强分布示意（定性，盘面下侧）
  \draw[red!50!black, thick] plot[smooth] coordinates {
    (-0.5,-0.9) (0.5,-0.85) (1.5,-0.8) (2.3,-0.78) (2.5,-1.3) (2.7,-0.82) (3.5,-0.85) (5,-0.9) (6,-0.92)
  };
  \node[annot, red!50!black] at (1.0,-1.15) {$p$分布};
  \node[annot, red!50!black] at (2.5,-1.5) {$\Delta p = T/A$};
  % 截面标记
  \node[annot] at (-0.1,0.9) {截面0};
  \node[annot] at (5.6,0.55) {截面2};
  % 流线
  \draw[stream, densely dotted] (0.2,0.3) .. controls (2.5,1.0) and (3.5,1.0) .. (5.5,0.5);
  \draw[stream, densely dotted] (0.2,-0.3) .. controls (2.5,-1.0) and (3.5,-1.0) .. (5.5,-0.5);
\end{tikzpicture}
\caption{桨盘动量理论控制体模型。理想桨盘为无限薄压强跳变面，流管截面在盘面处收缩。
盘前后速度连续但压强跳变$\Delta p = T/A$。诱导速度$v_i$使尾流速度达到$v_0+2v_i$。}
\label{fig:actuator_disk}
\end{figure}

取桨盘面积$A = \pi D^2 / 4$（$D$为桨径），定义远前方截面0（来流条件）、
桨盘截面和远后方尾流截面2构成流管控制体。假设流体无黏、不可压，
盘上压强跳变但速度连续（流量守恒要求截面相等则速度必连续），
盘前后压差在控制体轴向的积分为推力。

由质量守恒（流过各截面的质量流量相等）和轴向动量定理，
取控制体包含盘前和盘后全部流体：
\begin{equation}
    \dot{m} = \rho A v_p, \quad
    T = \dot{m} (v_2 - v_0)
    \label{eq:momentum_basic}
\end{equation}
其中$v_p$为盘面处流速，$v_0$为远前方来流速度（对于前飞机体即前行速度），
$v_2$为远后方尾流速度。

在盘面两侧分别应用伯努利方程（盘前：来流到盘面前；盘后：盘面后到尾流），
注意到盘面处速度连续但压强跳变$\Delta p = T/A$，可导出尾流速度为来流速度加两倍诱导速度：
$v_2 = v_0 + 2v_i$，其中诱导速度$v_i = v_p - v_0$为盘面处因桨盘做功而额外加速的流速。
代入式(\ref{eq:momentum_basic})即得推力的通用表达式：
\begin{equation}
    T = 2\rho A (v_0 + v_i) v_i
    \label{eq:thrust_general}
\end{equation}

这一表达式揭示了螺旋桨推力的基本物理：推力正比于空气质量流量（$\rho A v_p$）与速度增量（$v_2-v_0$）的乘积，
而速度增量恰为$2v_i$。增大推力有两种途径——增大桨盘面积（更高效地加速更多空气）或增大诱导速度（更快速地将空气向后抛）。
前者通过增大桨径实现，后者通过提高转速或增大桨距实现。

\textbf{悬停特例}（$v_0 = 0$）是分析推进系统基本性能的重要基准：
\begin{equation}
    T = 2\rho A v_i^2 \quad\Longrightarrow\quad
    v_i = \sqrt{\frac{T}{2\rho A}}
    \label{eq:hover_induced}
\end{equation}
悬停时的理想诱导功率$P_i = T v_i = T^{3/2} / \sqrt{2\rho A}$。
由此可得一个关键推论：在给定推力下，诱导功率与桨盘面积的关系为$P_i \propto 1/\sqrt{A} \propto 1/D$——
增大桨径是降低悬停功耗最有效的途径。这就是直升机旋翼大直径和多旋翼多桨盘设计的底层物理动机。
同时，理想悬停效率（figure of merit，$FM = P_i / P_{\text{actual}}$）是衡量
真实螺旋桨偏离理想动量理论程度的重要指标，在选型分析中应予以关注。

\textbf{前飞修正}（$v_0 > 0$）：在$v_0 \gg v_i$的前飞巡航工况，
诱导速度随前飞速度单调下降，推力公式趋近于$T \approx 2\rho A v_0 v_i$。
这意味着在同样的推力和桨盘面积下，前飞所需的诱导速度远小于悬停——换言之，
螺旋桨在前飞时比悬停时效率更高。这也解释了为什么固定翼飞机的巡航推进效率
远高于旋翼飞行器的悬停效率，以及为什么倾转旋翼构型在巡航模式下提供净收益。

对于拉力式（tractor）布局，前飞还带来一个额外效应：螺旋桨滑流加速流经机翼和机身的局部流速，
改变了当地动压分布，可能在一定程度上影响升力和力矩。在初步概念设计中，
滑流效应通常作为附加修正项处理，不在零阶动力学模型中展开。

\subsection{叶素理论：$k_T$、$k_Q$的物理来源}

动量理论将螺旋桨抽象为零厚度的压强跳变盘，不涉及桨叶的具体几何（弦长分布、桨距角、
翼型特性等）。叶素理论（blade element theory）\cite{mccormick1995aerodynamics}
弥补了这一不足：将桨叶分解为展向微元$dr$，每个微元视为以当地速度运动的二维翼型段，
分别计算其升力和阻力，然后沿展向积分得到总推力$T$和总扭矩$Q$。

对展向位置$r$处的微元，当地来流是旋转线速度$\omega r$（切向）与轴向速度
$v_0+v_i$（轴向，含诱导速度修正）的合成。微元产生的升力$dL$和阻力$dD$
按当地迎角（几何桨距角减流入角）由翼型气动特性给出。
将$dL$和$dD$分别投影到推力方向（轴向）和旋转平面（切向），沿桨叶展向积分，
再乘以桨叶数，得总推力和总扭矩。

将积分结果作量纲化处理（除以$\rho A (\omega R)^2$和$\rho A (\omega R)^2 R$等参考量），
引入无量纲推力系数$C_T$和扭矩系数$C_Q$：
\begin{equation}
    T = \frac{C_T \rho D^4}{4\pi^2} \omega^2 \equiv k_T \omega^2, \qquad
    Q = \frac{C_Q \rho D^5}{4\pi^2} \omega^2 \equiv k_Q \omega^2
    \label{eq:similarity}
\end{equation}
其中$C_T$、$C_Q$为无量纲系数，$D$为桨径，$\omega$为角速度。

式(\ref{eq:similarity})的形式揭示了螺旋桨性能的相似律。在固定桨距和固定前进比下，
$C_T$和$C_Q$近似为常数，推力与转速的平方成正比（$T \propto \omega^2$），
扭矩亦与转速的平方成正比（$Q \propto \omega^2$）。当桨距可调时，$C_T$和$C_Q$
均随桨距角近似线性增长——这意味着单位转速的推力和扭矩均可通过桨距调节来改变。

从量纲看：$k_T \propto \rho D^4$（推力系数$\times$密度$\times$特征面积），
$k_Q \propto \rho D^5$（扭矩系数$\times$密度$\times$特征面积$\times$特征长度）。
因此$k_Q / k_T \propto D$——反扭矩与推力之比正比于桨径。
这一量级关系对于评估滚转控制权限至关重要：差速反扭的滚转力矩量级为$k_Q \Delta(\omega^2) \sim k_Q \omega_0^2 \Delta\omega$，
而推力矢量力臂的俯仰/偏航力矩量级为$k_T \omega_0^2 \cdot a \cdot \delta$（$a$为力臂）。
两者之比约为$(k_Q / k_T) / a \sim D / a$，即桨径与力臂之比。

需要强调的是，$C_T$和$C_Q$并非真正的常数——它们随前进比$J = v_0/(nD)$（来流速度与桨尖速度之比）
和桨距角变化。在大包线（悬停$\to$高速巡航）过渡中，$C_T(J)$和$C_Q(J)$的变化可达$20\%$--$40\%$量级。
集总参数模型中的$k_T$和$k_Q$实质上是工作点处的局部线性化值，
这是该模型的主要近似来源之一。在概念级仿真中将$k_T$和$k_Q$取为固定的集总常数是工程上可接受的简化；
若需更高精度，可引入$J$相关的修正函数$k_T(J)$和$k_Q(J)$。

\subsection{反扭矩的物理机制与角动量守恒}

反扭矩（reaction torque）是螺旋桨驱动系统施加于机体的等大反向力矩。
深入理解反扭矩的物理来源对于设计差速滚转控制系统是必要的。

取转子和定子分别为研究对象的分析方法更为透彻。转子包括电机转子（永磁体）和螺旋桨组件，
对自身转轴的转动惯量为$J_p$，角速度为$\omega$。定子即固连于机体的电机定子（线圈）。
转子承受的外力矩有两个来源：定子通过电磁场施加的驱动力矩$\tau_e$（内部作用力，$+\tau_e$作用于转子），
以及空气作用于桨叶的气动阻力矩$Q$（外部力，$-Q$作用于转子，因为阻力矩方向与旋转方向相反）\cite{ducard2011fault}。

对转子列写角动量定理（转子自身转轴方向）：
\begin{equation}
    J_p \dot{\omega} = \tau_e - Q \quad\Longrightarrow\quad
    \tau_e = Q + J_p \dot{\omega}
    \label{eq:rotor_angular}
\end{equation}
其物理含义是：电磁驱动力矩必须同时克服气动阻力矩和转子加速所需的惯性力矩。

由牛顿第三定律，定子承受等大反向的力矩$-\tau_e$。从机体的视角看，
\textbf{机体承受的反扭矩方向与转子旋向相反}，大小为：
\begin{equation}
    \boxed{\tau_{\text{机体}} = \tau_e = k_Q \omega^2 + J_p \dot{\omega}}
    \label{eq:reaction_torque}
\end{equation}

此式揭示了反扭矩的两个组成部分及其物理性质：
\begin{itemize}
    \item \textbf{稳态反扭矩}$k_Q \omega^2$：在任何恒定转速下均存在，是螺旋桨将角动量持续传递给尾流所产生的反作用。
          方向始终与旋向相反，大小正比于$\omega^2$——转速加倍，反扭矩增至四倍；
    \item \textbf{惯性反扭矩}$J_p \dot{\omega}$：仅在转速变化时出现。加速阶段（$\dot{\omega} > 0$），
         惯性反扭矩叠加在稳态分量之上，机体瞬时承受更大的反扭；减速阶段（$\dot{\omega} < 0$），
         惯性反扭矩部分抵消气动阻力矩，机体承受的净反扭减小。
\end{itemize}

惯性反扭矩对飞行品质的影响不容忽视。在快速机动过程中，差速指令导致的转速剧烈变化
会通过惯性反扭矩在滚转通道产生瞬态扰动——因为差速本质上就是一台电机加速、另一台减速，
两者的$J_p \dot{\omega}$项叠加（加速电机惯性反扭增大，减速电机反扭减小），
在滚转轴上产生双倍的惯性反扭扰动。因此，\textbf{差速通道的带宽受限于电机时间常数和转子惯量}，
过快的差速变化不仅无法有效产生滚转力矩（电机来不及响应），还会引入额外的惯性扰动。

从更本质的角动量守恒视角看：将转子+排出流体+机体视为一个系统，三者角动量的变化率之和
在任何时刻精确为零。转子加速意味着转子角动量增大，同时排出流体的角动量流率$Q$也相应增大，
机体角动量（滚转速度）相应变化以维持总角动量守恒。这一视角有助于理解机动过程中的
角动量流动路径，特别是在分析快速滚转机动时的能量分配。

\subsection{电机一阶动态模型}

直流无刷电机（BLDC）的完整动态模型涉及电磁场、换相逻辑和功率电子器件的复杂交互。
对于飞行力学仿真和控制律设计，可以采用一阶惯性环节近似\cite{ducard2011fault}：
\begin{equation}
    \tau_m \dot{\omega} + \omega = \omega_{\text{cmd}}
    \label{eq:motor_lag}
\end{equation}
其中$\tau_m$为电机综合时间常数（电机+桨负载），$\omega_{\text{cmd}}$为目标转速指令。

$\tau_m$的物理构成由两部分决定：电机的电磁阻尼$k_t k_e / R$（$k_t$为扭矩常数，
$k_e$为反电动势常数，$R$为绕组电阻）提供基础的速度反馈；
螺旋桨的气动负载斜率$2 k_Q \omega_0$（随工作点增大）附加了额外的等效阻尼。
综合时间常数为：
\begin{equation}
    \tau_m = \frac{J_p}{k_t k_e / R + 2 k_Q \omega_0}
    \label{eq:tau_m}
\end{equation}

\textbf{关键物理洞察}：螺旋桨的平方负载（$Q \propto \omega^2$，其局部斜率为$2 k_Q \omega_0$）
本身就是天然的负反馈机制——当转速偏离工作点时，气动负载的变化倾向于将转速推回工作点。
这一机制在物理上等同于增大了等效阻尼系数，使$\tau_m$随$\omega_0$增大而减小（电机响应在高转速下更快）。

对于典型的小型无人机动力系统（$J_p \sim 2 \times 10^{-4}$\,kg$\cdot$m$^2$，
$k_t k_e / R \sim 2 \times 10^{-3}$\,N$\cdot$m$\cdot$s/rad），
螺旋桨负载的贡献$2 k_Q \omega_0$可与电磁阻尼同量级。
在高转速（高油门）下，$\tau_m$可降至低转速时的$60\%$--$70\%$——
这一效应在飞控增益调度中可通过油门-增益映射来近似补偿。

\subsection{旋转部件的陀螺效应}

旋转的电机-螺旋桨组件携带角动量$\bm{h} = J_p \omega \bm{n}$。
当机体有角速度$\bm{\Omega}$时，转子角动量矢量在惯性空间中的方向随体旋转而变化，
这一进动需要通过陀螺力矩来维持。由动系中的绝对导数公式（即矢量在不同参考系中时间导数的关系）：
\begin{equation}
    \left. \frac{d\bm{h}}{dt} \right|_{\text{惯性}}
    = \left. \frac{d\bm{h}}{dt} \right|_{\text{机体}}
    + \bm{\Omega} \times \bm{h}
    \label{eq:absolute_derivative}
\end{equation}

对于固连于机体的转子，在机体参考系中的角动量变化率$\dot{\bm{h}}|_{\text{机体}}$主要来自转速变化
（$J_p \dot{\omega} \bm{n}$）和摆座摆动导致的转轴方向变化。在大多数飞行动态中，
$\bm{\Omega} \times \bm{h}$项（进动效应）远大于$\dot{\bm{h}}|_{\text{机体}}$项
（转速变化的时间尺度通常比姿态变化慢）。忽略转子角动量大小的变化，得\textbf{陀螺力矩}：
\begin{equation}
    \boxed{\bm{M}_{\text{gyro}} = -\bm{\Omega} \times \bm{h}
           = -\bm{\Omega} \times (J_p \omega \bm{n})}
    \label{eq:gyro_moment}
\end{equation}

负号的含义是：陀螺力矩是机体角速度企图改变转子角动量方向时，转子对机体施加的反作用力矩——
本质上是对角动量方向改变的一种“抵抗”。

对于一个纵列双发构型，两台转子对机体的陀螺效应具有以下特征：
\begin{itemize}
    \item \textbf{机体俯仰角速率$q$}：前转子（转轴沿$+x_b$，$\bm{h}_f$沿$+x_b$）
          受$q$作用产生沿$+z_b$方向的进动趋势，陀螺力矩$-\bm{\Omega} \times \bm{h}_f$
          为$+y_b$方向力矩（偏航耦合）；尾转子（转轴沿$-x_b$，$\bm{h}_t$沿$-x_b$，
          但旋向相反导致$\bm{h}_t$亦沿$-x_b$）的进动产生$-y_b$方向力矩——
          \textbf{前后转子的俯仰-偏航陀螺耦合恰好反向，理想情况下相互抵消}。
          这正是反向旋转设计的一个常被忽视的重要优势；
    \item \textbf{机体偏航角速率$r$}：前转子受$r$作用产生$-y_b$方向陀螺力矩（俯仰耦合），
          尾转子产生$+y_b$方向陀螺力矩——再次相互抵消；
    \item \textbf{机体滚转角速率$p$}：前转子转轴沿$x_b$，$\bm{\Omega} \times \bm{h}$为零
          （角速度矢量与角动量矢量平行），不产生陀螺力矩。尾转子同理。
    \item \textbf{摆座摆动}：摆座绕$z_b$（前）或$y_b$（尾）的摆动本身是转轴的附加角速度
          $\bm{\Omega}_{\text{gimbal}}$，在摆动瞬间产生附加的陀螺力矩
          $\bm{M} = -\bm{\Omega}_{\text{gimbal}} \times \bm{h}$，
          量级约为$J_p \omega \dot{\delta}$。除非摆动速率极高（如伺服的阶跃响应），
          该项通常可以忽略。
\end{itemize}

\textbf{结论}：在纵列双发、反向旋转的构型中，陀螺效应的一阶项在前后转子之间相互抵消。
这一抵消性质是选择反向旋转的重要物理理由，远超出仅为了静态反扭矩对消的直觉。

% ===========================================================================
\section{推力矢量映射：从执行器到机体六维力/力矩}
\label{sec:mapping}
% ===========================================================================

推力矢量映射是将推进系统的推力与反扭矩（标量）转换为机体系六维力/力矩分量
（矢量）的桥梁。对于纵列双发正交单轴构型，这一映射具有特殊的简洁性，
其原因和后果是本节分析的核心。

\subsection{坐标系定义与符号约定}

在进行六维映射之前，必须明确定义各参考系及其转换关系：

\begin{table}[H]
\centering
\caption{坐标系定义与符号约定}
\begin{tabular}{ccl}
\toprule
符号 & 名称 & 定义 \\
\midrule
$\{\mathcal{I}\}$ & 惯性系 & 地面参考系，NED约定（$z_{\mathcal{I}}$向下），
                      原点可任意选取 \\
$\{\mathcal{B}\}$ & 机体系 & $x_b$指机头，$y_b$指右翼，$z_b$向下（右手系完备），
                      原点在飞行器瞬时质心 \\
$\{\mathcal{M}_f\}$ & 前电机系 & 与前电机固连，机身坐标系绕$z_b$旋转$\delta_f$后得到 \\
$\{\mathcal{M}_t\}$ & 尾电机系 & 与尾电机固连，机身坐标系绕$y_b$旋转$\delta_t$后得到 \\
\bottomrule
\end{tabular}
\end{table}

约定细节（确保符号的一致性）：
\begin{itemize}
    \item 前电机轴线沿$+x_b$方向（拉力式布局），正转（CW从后向前看），
          旋转矢量方向沿$+x_b$；
    \item 尾电机轴线沿$-x_b$方向（推进式布局），反转（CCW从后向前看），
          旋转矢量方向沿$+x_b$（反转意味着在尾电机自身参考系中为CCW，
          但从机体系看其角动量方向与推力方向相反）；
    \item 零摆角时两电机推力线均通过质心（CG）：$a$为前电机到CG的机体系$x$距离，
          前电机机体系位置$[+a, 0, 0]^T$；$b$为尾电机到CG的机体系$x$距离（正值），
          尾电机机体系位置$[-b, 0, 0]^T$；
    \item 摆角符号约定：前电机绕$+z_b$的右手旋转为正（$\delta_f > 0$表示推力向右偏）；
          尾电机绕$+y_b$的右手旋转为正（$\delta_t > 0$表示推力向下偏——在NED框架中即向下）。
\end{itemize}

\textbf{设计约束}：推力线过质心是保证通道解耦的必要几何条件。若这一条件不严格满足，
则零摆角推力本身会产生残余力矩，使解耦性退化。在实际飞行中，质心位置会随燃料消耗、
载荷变化而漂移，因此推力线设计应留有一定的不敏感区容差，或在SAS中引入前馈补偿。

\subsection{摆座运动学与推力方向}

单轴摆动即绕固定轴的旋转。尾电机绕$y_b$摆动$\delta_t$，
前电机绕$z_b$摆动$\delta_f$。对应的旋转矩阵为：
\begin{equation}
    \bm{R}_{y_b}(\delta_t) = \begin{bmatrix}
        \cos\delta_t & 0 & \sin\delta_t \\
        0 & 1 & 0 \\
        -\sin\delta_t & 0 & \cos\delta_t
    \end{bmatrix}, \qquad
    \bm{R}_{z_b}(\delta_f) = \begin{bmatrix}
        \cos\delta_f & -\sin\delta_f & 0 \\
        \sin\delta_f & \cos\delta_f & 0 \\
        0 & 0 & 1
    \end{bmatrix}
    \label{eq:gimbal_rotations}
\end{equation}

\textbf{前电机推力方向}（在$\{\mathcal{B}\}$系中的单位矢量）：
前电机自身推力方向在电机系中为$[1,0,0]^T$（拉力式向前），
经摆座旋转后在机体系中的方向为：
\begin{equation}
    \bm{u}_f = \bm{R}_{z_b}(\delta_f) \cdot [1, 0, 0]^T
    = [\cos\delta_f, \sin\delta_f, 0]^T
\end{equation}

\textbf{尾电机推力方向}（在$\{\mathcal{B}\}$系中的单位矢量）：
尾电机自身推力方向在电机系中为$-[1,0,0]^T$（推进式向后），
经摆座旋转后在机体系中的方向为：
\begin{equation}
    \bm{u}_t = \bm{R}_{y_b}(\delta_t) \cdot [-1, 0, 0]^T
    = [-\cos\delta_t, 0, \sin\delta_t]^T
\end{equation}

前电机转轴方向（角动量方向）与推力方向$\bm{u}_f$相同（正转CW，右手定则沿推力方向）。
尾电机转轴方向需要特别注意：尾电机是反转（CCW），其角动量方向与推力方向相反。
设尾电机推力方向为$\bm{u}_t$（即$[-\cos\delta_t, 0, \sin\delta_t]^T$），
则反转转子的角动量方向为$-\bm{u}_t$（即$[+\cos\delta_t, 0, -\sin\delta_t]^T$）。
这意味着尾电机的反扭矩在机体系中的方向为$+\bm{u}_t$的方向（反扭矩与转子旋向相反，
转子反转CCW，反扭矩方向为CW——即尾电机推力方向）。
在下面的力/力矩映射中将看到这一约定的影响。

\subsection{六维力/力矩映射的完整推导}

对两台电机分别计算推力矢量与反扭矩矢量的机体系分量，并按力臂生成力矩。
这是将标量执行器输出（$T_f, T_t, \tau_f, \tau_t$）转换为矢量力/力矩（$F_x, F_y, F_z, M_x, M_y, M_z$）
的核心变换。

\begin{figure}[H]
\centering
\begin{tikzpicture}[>=Stealth, scale=0.88,
  body/.style={draw=black!60, fill=black!6, rounded corners=2pt},
  arrow/.style={->, thick, >={Stealth[length=3mm]}},
  darrow/.style={->, thick, densely dashed, >={Stealth[length=2.5mm]}},
  annot/.style={font=\footnotesize\sffamily},
  dim/.style={<->, thin, draw=black!40},
]
  % 机体
  \draw[body] (-3.2,-1.6) rectangle (3.2,1.6);
  \fill (0,0) circle (2pt) node[annot, below right=2pt] {CG};
  % 坐标轴
  \draw[arrow] (-3.6,0) -- (3.6,0) node[annot, right] {$x_b$};
  \draw[arrow] (0,-2.0) -- (0,2.0) node[annot, above] {$y_b$};
  % z_b 垂直纸面向外
  \node[annot] at (0.4,0.3) {$\odot\,z_b$};
  % 前电机力矢量
  \draw[arrow, red!70!black, line width=1.2pt] (1.8,0) -- ++(35:2.0cm) node[annot, above] {$\bm{T}_f$};
  \draw[arrow, red!70!black, line width=1.2pt] (1.8,0) -- ++(0:2.0cm);
  \draw[red!70!black] (2.6,0) arc (0:35:8mm);
  \node[annot, red!70!black] at (2.95,0.35) {$\delta_f$};
  % 前电机位置
  \node[annot] at (1.8,-0.4) {前电机};
  \fill[red!30] (1.8,0) circle (3pt);
  % 前摆轴标注
  \node[annot, blue!60!black] at (1.8,-0.8) {摆轴$z_b$};
  % 尾电机力矢量
  \draw[arrow, blue!70!black, line width=1.2pt] (-1.6,0) -- ++(160:2.0cm) node[annot, above left] {$\bm{T}_t$};
  \draw[arrow, blue!70!black, line width=1.2pt] (-1.6,0) -- ++(180:2.0cm);
  \draw[blue!70!black] (-2.4,0) arc (180:160:8mm);
  \node[annot, blue!70!black] at (-2.7,0.25) {$\delta_t$};
  % 但尾电机在xy平面（俯视图）的力方向始终沿y=-tan(δ_t)方向...
  % 实际上尾电机绕y_b摆，在xz平面出分量，在xy平面（俯视图）看不到
  % 所以只显示沿-x_b基准方向
  \node[annot] at (-1.6,-0.4) {尾电机};
  \fill[blue!30] (-1.6,0) circle (3pt);
  % 力臂
  \draw[dim] (0,-1.2) -- node[annot, fill=white] {$a$} (1.8,-1.2);
  \draw[dim] (-1.6,-1.2) -- node[annot, fill=white] {$b$} (0,-1.2);
  % 力矩标注
  \draw[arrow, teal!60!black] (-0.3,1.5) arc (90:270:3mm and 6mm);
  \node[annot, teal!60!black] at (0.7,1.3) {$M_z = a T_f \sin\delta_f$};
  % 反扭矩示意
  \draw[darrow, orange!80!black] (1.8,0.5) arc (90:-180:4mm);
  \node[annot, orange!80!black] at (1.1,1.0) {$\tau_f$};
  \draw[darrow, orange!80!black] (-1.6,0.5) arc (90:270:4mm);
  \node[annot, orange!80!black] at (-2.5,1.0) {$\tau_t$};

  % ===== 侧视插图（小图） =====
  \begin{scope}[shift={(7.8,0)}]
    \draw[body] (-2.0,-1.2) rectangle (2.0,1.2);
    \fill (0,0) circle (2pt) node[annot, below right=2pt] {CG};
    \draw[arrow] (-2.4,0) -- (2.4,0) node[annot, right] {$x_b$};
    \draw[arrow] (0,1.6) -- (0,-1.6) node[annot, below] {$z_b$};
    \node[annot] at (0.4,-0.3) {$\otimes\,y_b$};
    % 尾电机在xz平面的力矢量
    \fill[blue!30] (-1.3,0) circle (3pt);
    \draw[arrow, blue!70!black, line width=1.2pt] (-1.3,0) -- ++(-20:2.0cm) node[annot, right=-1pt] {$\bm{T}_t$};
    \draw[arrow, blue!70!black, line width=1.2pt] (-1.3,0) -- ++(0:2.0cm);
    \draw[blue!70!black] (-0.5,0) arc (0:-20:8mm);
    \node[annot, blue!70!black] at (-0.2,-0.35) {$\delta_t$};
    % 前电机
    \fill[red!30] (1.5,0) circle (3pt);
    \draw[arrow, red!70!black, line width=1.2pt] (1.5,0) -- ++(0:1.8cm) node[annot, above] {$\bm{T}_f$};
    \node[annot] at (1.5,0.4) {前};
    \node[annot] at (-1.3,0.4) {尾};
    % 力臂
    \draw[dim] (0,0.9) -- node[annot, fill=white] {$a$} (1.5,0.9);
    \draw[dim] (-1.3,0.9) -- node[annot, fill=white] {$b$} (0,0.9);
    % 俯仰力矩
    \node[annot, teal!60!black] at (0.3,0.9) {$M_y = -b T_t \sin\delta_t$};
  \end{scope}

  \node[annot] at (3.25,-2.2) {(a) 俯视图：前摆产生偏航力矩};
  \node[annot] at (10.8,-2.2) {(b) 侧视图：尾摆产生俯仰力矩};
\end{tikzpicture}
\caption{推力矢量映射几何关系。前电机推力$\bm{T}_f$经摆角$\delta_f$产生侧向分量$T_f\sin\delta_f$，
经力臂$a$形成偏航力矩$M_z$；尾电机推力$\bm{T}_t$经摆角$\delta_t$产生垂向分量，经力臂$b$形成俯仰力矩$M_y$。
反扭矩$\tau_f,\tau_t$的差动控制滚转通道$M_x$。}
\label{fig:thrust_mapping}
\end{figure}

\textbf{前电机}（机体系位置$\bm{r}_f = [+a, 0, 0]^T$，拉力式，正转CW，反扭矩方向与旋向相反即沿$-\bm{u}_f$方向）：
\begin{align}
    \bm{F}_f &= T_f \begin{bmatrix} \cos\delta_f \\ \sin\delta_f \\ 0 \end{bmatrix},
    \quad
    \bm{\tau}_f = \tau_f \begin{bmatrix} \cos\delta_f \\ \sin\delta_f \\ 0 \end{bmatrix}
    \label{eq:front_force} \\
    \bm{M}_f &= \underbrace{\begin{bmatrix} +a \\ 0 \\ 0 \end{bmatrix} \times \bm{F}_f}_{\text{推力力臂}}
             - \underbrace{\bm{\tau}_f}_{\text{反扭矩}}
             = \begin{bmatrix}
                 -\tau_f \cos\delta_f \\
                 -\tau_f \sin\delta_f \\
                 a T_f \sin\delta_f
               \end{bmatrix}
    \label{eq:front_moment}
\end{align}
各分量的物理含义：
\begin{itemize}
    \item $M_{f,x} = -\tau_f \cos\delta_f$：前电机反扭矩在机体系$x$轴的分量，因反扭矩方向沿$-\bm{u}_f$，
          对滚转轴产生的是负贡献；
    \item $M_{f,y} = -\tau_f \sin\delta_f$：前摆角使反扭矩在$y$方向产生投影——这是偏航通道
          （前摆）向俯仰通道的一个二阶耦合；
    \item $M_{f,z} = a T_f \sin\delta_f$：推力$y$分量经力臂$a$产生偏航力矩——偏航主控力矩。
\end{itemize}

\textbf{尾电机}（机体系位置$\bm{r}_t = [-b, 0, 0]^T$，推进式，反转CCW，
反扭矩方向为推力方向即沿$\bm{u}_t$）：
\begin{align}
    \bm{F}_t &= T_t \begin{bmatrix} -\cos\delta_t \\ 0 \\ \sin\delta_t \end{bmatrix},
    \quad
    \bm{\tau}_t = \tau_t \begin{bmatrix} +\cos\delta_t \\ 0 \\ -\sin\delta_t \end{bmatrix}
    \label{eq:rear_force} \\
    \bm{M}_t &= \underbrace{\begin{bmatrix} -b \\ 0 \\ 0 \end{bmatrix} \times \bm{F}_t}_{\text{推力力臂}}
             + \underbrace{\bm{\tau}_t}_{\text{反扭矩}}
             = \begin{bmatrix}
                 +\tau_t \cos\delta_t \\
                 -b T_t \sin\delta_t \\
                 -\tau_t \sin\delta_t
               \end{bmatrix}
    \label{eq:rear_moment}
\end{align}
各分量的物理含义：
\begin{itemize}
    \item $M_{t,x} = +\tau_t \cos\delta_t$：尾电机反扭矩在机体系$x$轴的分量，因尾电机反转，
          其反扭矩方向与推力方向一致（沿$+\bm{u}_t$），对滚转轴贡献为正——与前电机的负贡献形成对消；
    \item $M_{t,y} = -b T_t \sin\delta_t$：推力$z$分量经力臂$b$产生俯仰力矩——俯仰主控力矩。
          注意负号：尾推力向下偏（$\delta_t>0$，$F_{t,z}>0$——NED中$z>0$为向下），
          产生抬头力矩（NED中$M_y<0$为机头向上）；
    \item $M_{t,z} = -\tau_t \sin\delta_t$：尾摆角使反扭矩在$z$方向产生投影——
          俯仰通道（尾摆）向偏航通道的一个二阶耦合。
\end{itemize}

合并前后电机贡献，得到\textbf{机体系六维推力矢量映射}（消去零项后）：
\begin{equation}
    \boxed{
    \begin{aligned}
        F_x &= T_f \cos\delta_f - T_t \cos\delta_t \\[2pt]
        F_y &= T_f \sin\delta_f \\[2pt]
        F_z &= T_t \sin\delta_t \\[6pt]
        M_x &= -\tau_f \cos\delta_f + \tau_t \cos\delta_t \\[2pt]
        M_y &= -b T_t \sin\delta_t - \tau_f \sin\delta_f \\[2pt]
        M_z &= a T_f \sin\delta_f - \tau_t \sin\delta_t
    \end{aligned}}
    \label{eq:six_component}
\end{equation}

\textbf{符号验证}（小角度展开确认物理直观正确性）：
\begin{itemize}
    \item $\delta_f > 0$（前摆右偏）：$F_y \approx T_f \delta_f > 0$（右侧力），
          $M_z \approx a T_f \delta_f > 0$（机头向右的偏航力矩），符合偏航直觉；
    \item $\delta_t > 0$（尾摆下偏）：$F_z \approx T_t \delta_t > 0$（NED向下力），
          $M_y \approx -b T_t \delta_t < 0$（抬头力矩），符合俯仰直觉；
    \item $\delta_f = \delta_t = 0$且$\tau_f = \tau_t$（等转速、零摆角、反扭对消）：
          $F_x = T_f - T_t$（净纵向推力），$F_y = F_z = 0$（无侧向/垂向力），
          $M_x = 0$（反扭对消，无滚转力矩），$M_y = M_z = 0$（无俯仰/偏航力矩）。
          在此理想状态下，所有控制力矩归零——这正是配平直飞的期望状态。
\end{itemize}

\textbf{耦合项分析}：式(\ref{eq:six_component})中包含两个交叉耦合项：
\begin{itemize}
    \item $M_y$中的$-\tau_f \sin\delta_f$：前摆角使反扭矩在$y$向产生投影，
          作用于俯仰轴，物理上是偏航摆动通过反扭矩投影对俯仰轴的耦合。
          量级估算：在典型参数下，该项约为$k_Q \omega_0^2 \cdot \delta_f$，
          而俯仰主控力矩为$b \cdot k_T \omega_0^2 \cdot \delta_t$。
          两者之比为$(k_Q / k_T b) \cdot (\delta_f/\delta_t) \sim (D/b) \cdot (\delta_f/\delta_t)$；
    \item $M_z$中的$-\tau_t \sin\delta_t$：尾摆角使反扭矩在$z$向产生投影，
          作用于偏航轴，物理上是俯仰摆动通过反扭矩投影对偏航轴的耦合。
          量级同理：为$k_Q \omega_0^2 \cdot \delta_t$，偏航主控力矩为$a \cdot k_T \omega_0^2 \cdot \delta_f$，
          两者之比为$(D/a) \cdot (\delta_t/\delta_f)$。
\end{itemize}

对于典型小型UAV参数（$D \approx 0.2$--$0.3$\,m，$a, b \approx 0.5$--$1.0$\,m），
耦合比在$20\%$--$60\%$范围——在工程上可接受但不完全可忽略。
当两通道同时大偏转时（如协调转弯），交叉耦合的叠加效应最显著，在SAS设计中宜通过
前馈补偿或耦合项直接建模来消除。

% ===========================================================================
\section{六自由度刚体动力学}
\label{sec:dynamics}
% ===========================================================================

整机作为刚体在三维空间中的运动由质心平动和绕质心的转动两部分组成。
本节建立完整的动力学和运动学方程，并讨论数值积分策略。

\subsection{牛顿-欧拉方程}

将飞行器视为刚体（弹性变形在概念设计阶段可忽略），其六自由度动力学由
牛顿第二定律（平动）和欧拉方程（转动）在机体系中联合描述\cite{stevens2015aircraft}。

\textbf{平动方程}。质心速度$\bm{v} = [u, v, w]^T$定义在机体系中。
然而牛顿第二定律$m\dot{\bm{v}}_{\mathcal{I}} = \sum \bm{F}$中的加速度是惯性系中的绝对加速度。
利用动系求导的坐标转换公式，机体系中观测的速度变化率加上因坐标系旋转带来的牵连项
等于惯性系绝对加速度在机体系的分量：
\begin{equation}
    m \left( \dot{\bm{v}} + \bm{\omega} \times \bm{v} \right)
    = \bm{F}_{\text{推力}} + \bm{F}_{\text{气动}} + m \bm{g}_b
    \label{eq:newton_body}
\end{equation}
其中$\bm{\omega} = [p, q, r]^T$为机体角速度。$\bm{\omega} \times \bm{v}$项
反映了即使$\bm{v}$在机体系中不变，因机体旋转导致$\bm{v}$在惯性系中的方向改变
所需的向心加速度。这是机体系中牛顿方程特有的附加项，在惯性系中写牛顿方程则不需要。

\textbf{转动方程}。对机体系原点（质心）取矩，并考虑旋转部件（电机+螺旋桨）
的角动量效应：
\begin{equation}
    \bm{I} \dot{\bm{\omega}} + \bm{\omega} \times (\bm{I} \bm{\omega})
    + \bm{\omega} \times \bm{h}_{\text{转子}}
    = \bm{M}_{\text{推力}} + \bm{M}_{\text{气动}}
    \label{eq:euler_body}
\end{equation}
各物理项的含义依次为：
\begin{itemize}
    \item $\bm{I} \dot{\bm{\omega}}$：机体角加速度引起的惯性力矩（转动惯量乘角加速度）。
          对于主轴坐标系，$\bm{I} = \operatorname{diag}(I_x, I_y, I_z)$；
    \item $\bm{\omega} \times (\bm{I} \bm{\omega})$：因机体角速度矢量与角动量矢量
          不共线（除非绕主轴旋转）而产生的陀螺进动力矩。这一项引起三轴之间的动态耦合——
          例如滚转（$p$）和偏航（$r$）均不为零时，机身绕$y$轴产生$(I_x-I_z)pr$的惯性耦合力矩；
    \item $\bm{\omega} \times \bm{h}_{\text{转子}}$：转子总角动量
          $\bm{h}_{\text{转子}} = \sum_i J_{p,i} \omega_i \bm{n}_i$
          因机体旋转产生的陀螺力矩（已在第\ref{sec:propulsion}节详细讨论）。
          注意此处是$+\bm{\omega} \times \bm{h}$（式\ref{eq:euler_body}左侧），
          与式(\ref{eq:gyro_moment})的$-\bm{\Omega} \times \bm{h}$看似符号相反，
          实则为同一物理量在方程两侧的不同表述——左侧是因，右侧是果。
\end{itemize}

\textbf{重力在机体系的分量}。重力在惯性系中恒为$[0, 0, mg]^T$（NED约定，$z_{\mathcal{I}}$向下）。
通过姿态旋转矩阵的逆变换（从惯性系到机体系）得：
\begin{equation}
    \bm{g}_b = \bm{R}_{\mathcal{I}\mathcal{B}}
    \begin{bmatrix} 0 \\ 0 \\ mg \end{bmatrix}
    = \bm{R}_{\mathcal{B}\mathcal{I}}^T
    \begin{bmatrix} 0 \\ 0 \\ mg \end{bmatrix}
    \label{eq:gravity_body}
\end{equation}

\textbf{气动力建模}。在概念级仿真中，气动力可用简化的集总参数模型\cite{beard2012small}：
动压$q_\infty = \frac{1}{2}\rho V^2$（$V$为空速），迎角$\alpha = \arctan(w/u)$，
侧滑角$\beta = \arcsin(v/V)$。气动力的完整建模需要气动导数数据库（通常来自CFD或风洞试验），
不在纯推力矢量分析的范围内。在推力矢量构型的初步分析中，可先忽略气动力以聚焦推力矢量的
控制特性，随后逐步引入气动项以评估气动-推力矢量的交互效应。

\subsection{姿态运动学：四元数表示}

采用四元数$\bm{q} = [q_0, q_1, q_2, q_3]^T$（标量在前约定）
描述机体到惯性系的旋转。选择四元数而非欧拉角的主要理由是在$\theta = \pm 90^\circ$处
欧拉角存在万向锁奇异——3-2-1欧拉角序列在$\theta = \pm 90^\circ$时，
$\phi$和$\psi$描述同一物理旋转自由度，导致运动学方程退化。
四元数无此奇异点，对于可能包含大俯仰角的推力矢量构型尤其必要\cite{markley2014fundamentals,kuipers1999quaternions}。

姿态运动学方程（机体系角速度$\bm{\omega}_b = [p, q, r]^T$驱动四元数演化）：
\begin{equation}
    \dot{\bm{q}} = \frac{1}{2} \bm{q} \otimes
    \begin{bmatrix} 0 \\ p \\ q \\ r \end{bmatrix}
    \label{eq:quaternion_kinematics}
\end{equation}
写成矩阵形式（便于数值实现）：
\begin{equation}
    \begin{bmatrix} \dot{q}_0 \\ \dot{q}_1 \\ \dot{q}_2 \\ \dot{q}_3 \end{bmatrix}
    = \frac{1}{2}
    \begin{bmatrix}
        0 & -p & -q & -r \\
        p &  0 &  r & -q \\
        q & -r &  0 &  p \\
        r &  q & -p &  0
    \end{bmatrix}
    \begin{bmatrix} q_0 \\ q_1 \\ q_2 \\ q_3 \end{bmatrix}
    \label{eq:quaternion_matrix}
\end{equation}

四元数到方向余弦矩阵（$\{\mathcal{B}\}\to\{\mathcal{I}\}$）的转换：
\begin{equation}
    \bm{R}_{\mathcal{B}\mathcal{I}} =
    \begin{bmatrix}
        1-2(q_2^2+q_3^2) & 2(q_1 q_2 - q_0 q_3) & 2(q_1 q_3 + q_0 q_2) \\
        2(q_1 q_2 + q_0 q_3) & 1-2(q_1^2+q_3^2) & 2(q_2 q_3 - q_0 q_1) \\
        2(q_1 q_3 - q_0 q_2) & 2(q_2 q_3 + q_0 q_1) & 1-2(q_1^2+q_2^2)
    \end{bmatrix}
    \label{eq:dcm_from_quaternion}
\end{equation}

3-2-1欧拉角（偏航-俯仰-滚转，用于显示和人机界面）从DCM提取：
\begin{equation}
    \begin{aligned}
        \phi &= \operatorname{atan2}(R_{23}, R_{33}) \\
        \theta &= -\arcsin(R_{13}) \\
        \psi &= \operatorname{atan2}(R_{12}, R_{11})
    \end{aligned}
    \label{eq:euler_from_dcm}
\end{equation}

\textbf{四元数的数值维护}：每次积分步后，四元数因浮点舍入误差的积累会逐渐偏离单位模长
（$\|\bm{q}\| \neq 1$）。若不加处理，范数漂移会导致DCM失去正交性，
进而使坐标变换产生系统性畸变。解决方案是对每一积分子步后的四元数执行归一化：
$\bm{q} \leftarrow \bm{q} / \|\bm{q}\|$。这是飞行仿真中的标准做法。

\subsection{惯性系位置与速度}

机体系速度经姿态变换到惯性系后积分得到惯性系位置：
\begin{equation}
    \dot{\bm{p}}_{\mathcal{I}} = \bm{R}_{\mathcal{B}\mathcal{I}} \bm{v}_{\mathcal{B}}
    \label{eq:position_integration}
\end{equation}

\subsection{子步调度与数值稳定性}

对于显式欧拉积分方法，数值稳定性受限于系统最快时间常数的两倍（Nyquist-Shannon型约束
在数值ODE中的对应物是$\Delta t < 2/\lambda_{\max}$，其中$\lambda_{\max}$为系统特征值的最大实部绝对值）。
在六自由度飞行仿真中，最快的时间常数通常来自气动阻尼通道或高增益控制回路。
以滚转通道为例，气动滚转阻尼的时间常数为$\tau_p \approx I_x / (q_\infty S b^2 C_{l_p} / (2V))$，
在典型参数下可能小至$10$--$20$\,ms。因此要求积分步长$\Delta t \lesssim 4$\,ms以保证
数值稳定性。

\textbf{子步调度策略}：主循环（渲染+物理）以屏幕刷新率（通常16--17\,ms）或固定频率运行，
物理积分采用更小的子步。设主循环步长为$\Delta t$，取$n = \lceil \Delta t / \Delta t_{\max} \rceil$
个子步（$\Delta t_{\max}$通常取4\,ms），每子步以$h = \Delta t / n$逐步积分。
子步内的计算流程为：
\begin{enumerate}[label=(\arabic*)]
    \item 推进系统：根据当前油门与差速指令计算推力$T$与反扭矩$\tau$；
    \item 推力矢量映射：根据当前摆角计算六维力/力矩分量；
    \item 欧拉方程积分：$\bm{\omega} \leftarrow \bm{\omega} + h \cdot \dot{\bm{\omega}}$；
    \item 四元数更新：$\bm{q} \leftarrow \bm{q} + h \cdot \dot{\bm{q}}$，随后归一化；
    \item 平动速度更新：$\bm{v} \leftarrow \bm{v} + h \cdot \dot{\bm{v}}$；
    \item 惯性系位置更新：$\bm{p}_{\mathcal{I}} \leftarrow \bm{p}_{\mathcal{I}} + h \cdot \dot{\bm{p}}_{\mathcal{I}}$。
\end{enumerate}
主循环步长上限（frame cap）设为$50$\,ms，防止因浏览器标签切换、后台暂停等导致的
大时间跳跃引起的积分发散。

% ===========================================================================
\section{控制分配：通道解耦的几何机理}
\label{sec:allocation}
% ===========================================================================

控制分配（control allocation）是连接控制律输出（期望力矩）与执行器输入（摆角/差速指令）
的数学环节。本构型在控制分配方面的最显著特征是：正交摆轴的几何设计使效能矩阵退化为对角阵，
控制分配从欠定优化问题简化为三通道直接映射。

\subsection{小角度近似与线性化效能矩阵}

在小摆角近似下（$\sin\delta \approx \delta$，$\cos\delta \approx 1$，$\delta \ll 1$），
并假设前后推力近似相等（$T_f \approx T_t \approx T_0$），式(\ref{eq:six_component})
中的力矩方程线性化为：
\begin{equation}
    \begin{bmatrix} M_x \\ M_y \\ M_z \end{bmatrix}
    \approx
    \underbrace{
    \begin{bmatrix}
        -k_Q \omega_f^2 + k_Q \omega_t^2 \\
        -b T_t \delta_t \\
        a T_f \delta_f
    \end{bmatrix}
    }_{\text{主控通道}}
    +
    \underbrace{
    \begin{bmatrix}
        0 \\
        -\tau_f \delta_f \\
        -\tau_t \delta_t
    \end{bmatrix}
    }_{\text{交叉耦合（二阶）}}
    \label{eq:moment_linearized}
\end{equation}

\textbf{差速变量的参数化}：定义差速变量$\Delta\omega$，使两电机目标转速为：
\begin{equation}
    \omega_f^{\text{target}} = \omega_0 \sqrt{1 + \Delta\omega}, \quad
    \omega_t^{\text{target}} = \omega_0 \sqrt{1 - \Delta\omega}
    \label{eq:diff_def}
\end{equation}
其中$\omega_0$为基准转速（由油门决定），$\Delta\omega \in [-1, 1]$为归一化差速指令。
此分配方案的物理精妙之处在于：$\omega_f^2 + \omega_t^2 = 2\omega_0^2$严格不变——
\textbf{在任何差速指令下，两电机转速平方之和为常数}。由$T \propto \omega^2$，
一阶近似下总推力保持不变，差速仅影响扭矩差而不影响总推力。
这一性质使操纵手的油门-滚转解耦直观感受得以保证：推油门只变化总推力，
滚转指令只变化差速，两者在推力一阶近似下独立。

令基准反扭矩$\tau_0 = k_Q \omega_0^2$，滚转力矩的差速分量为：
\begin{equation}
    \Delta M_x^{\text{差速}} = k_Q (\omega_t^2 - \omega_f^2)
    = k_Q (\omega_0^2(1-\Delta\omega) - \omega_0^2(1+\Delta\omega))
    = -2 k_Q \omega_0^2 \Delta\omega
    = -2 \tau_0 \Delta\omega
    \label{eq:differential_roll}
\end{equation}

合并，得\textbf{线性化控制效能矩阵}：
\begin{equation}
    \boxed{
    \begin{bmatrix} M_x \\ M_y \\ M_z \end{bmatrix}
    =
    \underbrace{
    \begin{bmatrix}
        -2\tau_0 & 0 & 0 \\
        0 & -b T_0 & 0 \\
        0 & 0 & a T_0
    \end{bmatrix}
    }_{\bm{B}(\omega_0)}
    \begin{bmatrix} \Delta\omega \\ \delta_t \\ \delta_f \end{bmatrix}
    +
    \text{耦合项} + \text{气动项}
    }
    \label{eq:effectiveness_matrix}
\end{equation}

\subsection{对角化效能矩阵的几何根源}

效能矩阵$\bm{B}(\omega_0)$为严格对角阵——这是构型几何设计而非数值近似的结果。
对角线外的零项来自于以下几何约束的精确满足：

\begin{itemize}
    \item \textbf{前电机绕$z_b$轴摆动}：推力方向的任何变化均限制在$xy$平面内。
          推力$y$分量$T_f \sin\delta_f$经力臂$a$产生偏航力矩$M_z = a T_f \sin\delta_f$。
          由于摆轴为$z_b$，前电机推力不可能产生$z$向分量，因此不可能产生俯仰力矩（$M_y$）
          ——除了反扭矩投影的二阶耦合$-\tau_f\sin\delta_f$；
    \item \textbf{尾电机绕$y_b$轴摆动}：推力方向的任何变化均限制在$xz$平面内。
          推力$z$分量$T_t\sin\delta_t$经力臂$b$产生俯仰力矩$M_y = -b T_t\sin\delta_t$。
          由于摆轴为$y_b$，尾电机推力不可能产生$y$向分量，因此不可能产生偏航力矩（$M_z$）
          ——除了反扭矩投影的二阶耦合$-\tau_t\sin\delta_t$；
    \item \textbf{差速通道}：仅影响反扭矩$M_x$，在$M_y$和$M_z$中的投影为零
          （反扭矩矢量的$y$和$z$分量分别来自前/尾电机的摆角投影，差速本身不改变摆角）。
\end{itemize}

这一设计使控制分配问题从通用的$\bm{u} = \bm{B}^+ \bm{M}_{\text{cmd}}$（其中
$\bm{B}^+$为$3 \times 3n$的伪逆，$3n > 3$，需在线求解）
退化为简单的\textbf{三通道独立映射}：
\begin{equation}
    \Delta\omega = -\frac{M_x^{\text{cmd}}}{2\tau_0}, \quad
    \delta_t = -\frac{M_y^{\text{cmd}}}{b T_0}, \quad
    \delta_f = \frac{M_z^{\text{cmd}}}{a T_0}
    \label{eq:direct_mapping}
\end{equation}

这是本构型相对于通用矢量推力平台的根本优势——
用几何设计换取了算法的简洁性\cite{johansen2013control,oppenheimer2010control}。

\subsection{耦合项量级评估与补偿策略}

线性化后的交叉耦合项来源于反扭矩在非主控轴方向的投影。一阶（小摆角）近似下：
\begin{align}
    \Delta M_y^{\text{耦合}} &= -\tau_f \delta_f
    \approx -k_Q \omega_f^2 \cdot \delta_f
    \label{eq:coupling_pitch} \\
    \Delta M_z^{\text{耦合}} &= -\tau_t \delta_t
    \approx -k_Q \omega_t^2 \cdot \delta_t
    \label{eq:coupling_yaw}
\end{align}

\textbf{耦合-主控比}的量级评估。以俯仰通道为例，在典型巡航工况下，
主控力矩量级为$b T_f \delta_t \approx b \cdot k_T \omega_0^2 \cdot \delta_t$，
耦合力矩量级为$\tau_f \delta_f \approx k_Q \omega_0^2 \cdot \delta_f$。在俯仰和偏航
同时大幅偏转的极限工况，耦合比约为：
\begin{equation}
    \frac{\text{耦合}}{\text{主控}}
    = \frac{k_Q \omega_0^2 \delta_f}{b \cdot k_T \omega_0^2 \delta_t}
    = \frac{k_Q}{k_T b} \cdot \frac{\delta_f}{\delta_t}
    \label{eq:coupling_ratio}
\end{equation}

由量纲分析$k_Q / (k_T b) \sim D / b$。对于典型小型UAV
（$D \approx 0.2$--$0.3$\,m，$b \approx 0.5$--$1.0$\,m），
这一比值在$0.2$--$0.6$范围。在正常机动下$\delta_f$和$\delta_t$通常同量级
（耦合比为$20\%$--$60\%$），在极限情况下若一通道偏转显著大于另一通道，
耦合比可能接近甚至超过$100\%$。

\textbf{工程对策}：
\begin{enumerate}[label=(\arabic*)]
    \item \textbf{显式前馈补偿}：在控制律中直接计算耦合项并前馈抵消。
          由$\Delta \delta_t^{\text{comp}} = +(k_Q \omega_f^2 / (b T_0)) \cdot \delta_f$，
          在计算俯仰摆角指令时加入此项以抵消偏航摆动的耦合效应；
    \item \textbf{SAS鲁棒性}：设计足够高的SAS反馈增益，使闭环系统对耦合扰动具有足够衰减。
          耦合扰动在闭环下被抑制了$1/(1+GH)$倍，其中$GH$为回路增益；
    \item \textbf{非线性精确映射}：当计算资源允许时，直接使用非线性映射（式\ref{eq:six_component}）
          配合牛顿迭代求解$\delta_f, \delta_t, \Delta\omega$，精确消除耦合项的影响。
          此方法在摆角$>25^\circ$的大机动中是必要的，因为$\sin\delta \approx \delta$
          的近似精度严重退化（$25^\circ$时$\sin\delta/\delta \approx 0.95$，
          误差约$5\%$；$45^\circ$时误差已达$10\%$）。
\end{enumerate}

\subsection{低油门滚转控制退化}

差速滚转效能$\partial M_x / \partial \Delta\omega = -2 k_Q \omega_0^2 \propto T_0$——
\textbf{滚转控制力矩的大小与总推力成正比}。这是差速反扭滚转机制的物理固有特性，
无法通过控制律增益来补偿（因为即使增益无穷大，执行器输出已饱和）。

\textbf{物理极限}：
\begin{itemize}
    \item 油门$100\%$：滚转效能最大，机动性最好；
    \item 油门$50\%$：滚转效能降至$25\%$（因$\propto \omega_0^2 \propto \text{thr}$，实际关系接近$T_0 \propto \text{thr}^{2/3}$量级，但比线性退化更快）；
    \item 油门$\to 0$：滚转控制失效——在此极限下，从推进系统看飞行器退化为无控刚体。
\end{itemize}

这一特性对飞行安全有直接意义：进近着陆段通常为低油门状态，而此时恰恰是滚转控制
最需要的阶段（侧风修正、对准跑道中心线）。缓解措施包括：
保持进近油门不低于某个安全阈值（即使这意味着需要更陡的下滑角），
或引入气动副翼作为低油门的滚转控制补充（但这增加了构型复杂度）。

\subsection{与通用过驱动推力矢量平台的对比}

通用多推进器矢量推力平台的设计空间可以是过驱动（over-actuated）或全驱动（fully-actuated）。
对于$n$个推进器、每个推进器具有2自由度万向节的情况，执行空间维数为$3n$
（$n$个推进器$\times$ 3维力矢量），映射到6维力/力矩空间。
当$3n > 6$（即$n \ge 3$）时形成过驱动系统\cite{ding2018modeling}。

过驱动系统的控制分配是一个经典的欠定问题：
给定期望力矩$\bm{M}_{\text{cmd}} \in \mathbb{R}^3$（或期望六维力/力矩$\in \mathbb{R}^6$），
求执行器输入$\bm{u} \in \mathbb{R}^{3n}$，使$\bm{B}\bm{u} = \bm{M}_{\text{cmd}}$。
由于$3n > 3$（或$>6$），方程有无穷多解。需要在解空间中选取“最优”解，
通常通过求解以下优化问题之一\cite{johansen2013control}：
\begin{itemize}
    \item \textbf{最小二范数（伪逆）}：$\min_{\bm{u}} \|\bm{u}\|^2$ s.t. $\bm{B}\bm{u} = \bm{M}_{\text{cmd}}$。
          解为$\bm{u} = \bm{B}^T(\bm{B}\bm{B}^T)^{-1}\bm{M}_{\text{cmd}}$。简单但未考虑执行器限幅；
    \item \textbf{加权伪逆}：$\min_{\bm{u}} \bm{u}^T\bm{W}\bm{u}$ s.t. $\bm{B}\bm{u} = \bm{M}_{\text{cmd}}$。
          通过权重矩阵$\bm{W}$引入不同执行器优先级（如优先使用大推力电机、保护接近饱和的执行器）；
    \item \textbf{二次规划}：$\min_{\bm{u}} \|\bm{W}_u(\bm{u} - \bm{u}_p)\|^2$ s.t.
          $\bm{B}\bm{u} = \bm{M}_{\text{cmd}}$且$\bm{u}_{\min} \le \bm{u} \le \bm{u}_{\max}$。
          在满足等式和不等式约束的前提下最小化与标称输入的偏差。
\end{itemize}

在线求解这些优化问题需要迭代算法（如active-set法或内点法），
在计算资源受限的嵌入式飞控平台上是实质性挑战。本构型通过将执行空间维度
从$3n$降至3（恰等于控制目标维度），使$\bm{B}$方阵且对角——
\textbf{直接用几何设计换取算法的零计算开销}。这是构型设计的核心价值主张。

% ===========================================================================
\section{线性化与稳定性模态分析}
\label{sec:stability}
% ===========================================================================

线性化是分析飞行器稳定性特性的经典方法。在本节中，我们将非线性六自由度动力学方程
在工作点附近线性化，得到状态空间模型，并分析纵向和横航向的各个特征模态。

\subsection{小扰动线性化方法}

取配平状态为定直平飞（稳态、水平、无侧滑，$\phi_0 = \beta_0 = 0$，
$p_0 = q_0 = r_0 = 0$）。总状态向量为（采用分离形式以适应经典模态分析的惯例）：
\begin{align}
    \bm{x}_{\text{long}} &= [\Delta u, \Delta w, \Delta q, \Delta \theta]^T \quad \text{（纵向，4阶）} \\
    \bm{x}_{\text{lat}} &= [\Delta v, \Delta p, \Delta r, \Delta \phi]^T \quad \text{（横航向，4阶）}
\end{align}

在配平点处对非线性函数$\bm{f}(\bm{x}, \bm{u})$作Taylor展开，
保留一阶项得到线性状态空间模型（形式上的A/B矩阵）：
\begin{equation}
    \Delta \dot{\bm{x}} = \underbrace{\frac{\partial \bm{f}}{\partial \bm{x}}\bigg|_0}_{\bm{A}}
    \Delta \bm{x} + \underbrace{\frac{\partial \bm{f}}{\partial \bm{u}}\bigg|_0}_{\bm{B}}
    \Delta \bm{u}
    \label{eq:linear_state_space}
\end{equation}
其中控制输入$\Delta \bm{u} = [\Delta \delta_t, \Delta \delta_f, \Delta \omega]^T$。
偏导数矩阵的每个元素（稳定导数与控制导数）在概念设计阶段通过有限差分或解析求导获得。
在高可信度阶段，这些导数应由CFD计算或风洞试验标定。

\subsection{纵向模态的物理本质}

纵向四阶系统（$u, w, q, \theta$）的A矩阵具有如下结构：
\begin{equation}
    \bm{A}_{\text{long}} =
    \begin{bmatrix}
        X_u & X_w & X_q - w_0 & -g\cos\theta_0 \\
        Z_u & Z_w & Z_q + u_0 & -g\sin\theta_0 \\
        M_u & M_w & M_q & 0 \\
        0 & 0 & 1 & 0
    \end{bmatrix}
    \label{eq:longitudinal}
\end{equation}
其中各导数的物理含义为：$X_u = (1/m)\partial F_x/\partial u$（速度导数，通常为负——速度增大$\to$阻力增大），
$Z_w = (1/m)\partial F_z/\partial w$（法向力对$w$的导数，决定升力线斜率），
$M_q = (1/I_y)\partial M_y/\partial q$（俯仰阻尼导数，恒为负），
$M_w = (1/I_y)\partial M_y/\partial w$（静安定性导数，$M_w < 0$表示静稳定）。

纵向系统通常呈现经典的两模态结构，这是飞行力学教科书中的标准结论\cite{etkin1982dynamics,stengel2004flight}：

\textbf{短周期模态}（short period mode）：以$w$（迎角）和$q$（俯仰速率）
为主变量的快速、高阻尼振荡（或两个实根）。典型周期为1--3秒。其近似特征频率
和阻尼由静稳定导数$M_w$（$C_{m\alpha}$）和俯仰阻尼$M_q$共同决定：
\begin{equation}
    \omega_{sp} \approx \sqrt{-M_w \cdot Z_q/u_0}, \quad
    \zeta_{sp} \approx \frac{-(M_q + M_{\dot{w}} \cdot Z_q)}{2\omega_{sp}}
\end{equation}

对于纵列双发构型，在SAS闭环下，尾摆的附加俯仰力矩（$\partial M_y / \partial \delta_t = -b T_0$，
经SAS增益$k_q$转换为有效俯仰阻尼）进一步加速了短周期收敛。
这相当于在气动阻尼$M_q$的基础上附加了推力矢量阻尼$M_q^{\text{TV}} = k_q \cdot (-b T_0) / I_y$，
使闭环短周期阻尼显著高于开环阻尼。

\textbf{长周期模态}（phugoid mode）：以$u$（速度）和$\theta$（俯仰角）为主变量的
缓慢、低阻尼振荡。周期通常在20--60秒。其物理本质是飞行器动能（$\sim mu^2/2$）
和势能（$mgh$）之间的缓慢交换：飞行器抬头$\to$速度下降$\to$升力减小$\to$低头
$\to$速度恢复$\to$升力恢复$\to$再次抬头。推力矢量的摆角变化对长周期模态的影响
通常不显著——长周期主要取决于总能量交换（升阻比和推力-阻力平衡），
而非快速力矩平衡，而推力矢量主要影响后者。

\textbf{推力矢量对纵向稳定性的贡献}：推力矢量的静稳定效应可通过等效静稳定裕度
来量化。定义$\partial C_m / \partial \alpha$为全机俯仰力矩系数对迎角的导数。
对于推力矢量构型，该导数包含气动贡献$C_{m\alpha}^{\text{aero}}$和
推力矢量贡献$C_{m\alpha}^{\text{TV}}$，后者来源于推力矢量随迎角变化导致的
力臂效应变化和来流-转速耦合（前进比变化导致的推力变化）。在初步分析中，
推力矢量的等效静稳定贡献通常为正（增强稳定性），因为低头力矩随迎角增大而增大
（尾摆推力的法向分量经力臂产生的低头力矩增大）。

\subsection{横航向模态的物理本质}

横航向四阶系统（$v, p, r, \phi$）的A矩阵：
\begin{equation}
    \bm{A}_{\text{lat}} =
    \begin{bmatrix}
        Y_v & Y_p + w_0 & Y_r - u_0 & g\cos\theta_0 \\
        L_v & L_p & L_r & 0 \\
        N_v & N_p & N_r & 0 \\
        0 & 1 & \tan\theta_0 & 0
    \end{bmatrix}
    \label{eq:lateral}
\end{equation}
各关键导数：$L_v = (1/I_x)\partial M_x/\partial v$（上反角效应，$L_v < 0$表示正上反角——有利的滚转稳定性），
$L_p = (1/I_x)\partial M_x/\partial p$（滚转阻尼，恒为负），
$N_v = (1/I_z)\partial M_z/\partial v$（航向静稳定性，$N_v > 0$为稳定——侧滑产生顺风向偏航力矩），
$N_r = (1/I_z)\partial M_z/\partial r$（偏航阻尼，恒为负）。

横航向系统呈现三个特征模态（通常为一个快衰减实根、一对共轭复根和一个慢实根）\cite{etkin1982dynamics}：

\textbf{滚转收敛模态}（roll subsidence）：以$p$（滚转速率）为绝对主导的单实根模态。
近似特征值$\lambda_R \approx L_p < 0$。无控状态下，气动滚转阻尼$C_{l_p}$提供收敛，
典型时间常数约为$0.5$--$2$秒。纵列双发构型在SAS闭环下，差速反扭提供附加的人工滚转阻尼
$L_p^{\text{SAS}} = k_p \cdot (-2 k_Q \omega_0^2) / I_x $，进一步缩短收敛时间。

\textbf{荷兰滚模态}（Dutch roll）：$v$（侧滑角）和$r$（偏航速率）的耦合振荡。
在滚转自由度相对不活跃时（小上反角或无滚转输入），荷兰滚近似为侧滑和偏航的
二阶振荡系统。其近似特征频率由航向静稳定性$N_v$（$C_{n\beta}$）决定：
\begin{equation}
    \omega_{DR} \approx \sqrt{N_v \cdot u_0 / I_z}
\end{equation}
荷兰滚的阻尼来源包括偏航阻尼$N_r$和侧力导数$Y_v$的间接贡献。
一个关键设计准则是：荷兰滚的阻尼必须足够大（通常$\zeta_{DR} > 0.05$--$0.1$），
否则飞行器在湍流中的偏航-滚转耦合响应将严重影响乘客舒适度（对于载人飞行器）
或航向保持精度（对于无人机）。

对于纵列双发构型，前摆的偏航力矩$\partial M_z / \partial \delta_f = a T_0$
经SAS偏航阻尼增益$k_r$转换为有效偏航阻尼$N_r^{\text{SAS}} = -k_r \cdot a T_0 / I_z < 0$，
显著增强了荷兰滚的阻尼。实际上，推力矢量SAS是荷兰滚阻尼的最有效增强手段之一——
它直接作用于偏航轴，产生与偏航速率成比例的对抗力矩。

\textbf{螺旋模态}（spiral mode）：$\phi$（滚转角）和$r$（偏航速率）的缓慢耦合发散或收敛。
近似特征值（在荷兰滚频率远大于螺旋模态频率的假设下）：
\begin{equation}
    \lambda_S \approx \frac{g}{u_0} \cdot
    \frac{L_v N_r - N_v L_r}{L_v N_p - N_v L_p}
\end{equation}
螺旋模态的物理：一侧滚转（$\phi$偏转）$\to$升力倾斜$\to$侧向力$\to$侧滑
$\to$航向静稳定性产生顺风向偏航力矩$\to$偏航加剧了滚转方向（在协调转弯中，
偏航的外翼速度高于内翼，产生更大的升力，进一步增大滚转角）。

螺旋模态可能随时间缓慢发散（典型的倍幅时间达数十秒至数分钟），
对飞行员而言是容易觉察和纠正的——与高频率的荷兰滚/短周期发散相比，
缓慢发散在实际飞行中通常被认为是可接受的。但在自动驾驶仪和导航控制中，
未抑制的螺旋发散会导致持续的方向漂移，必须通过航向保持控制律来抑制。

\subsection{SAS对模态的影响}

SAS闭环改变了系统矩阵$\bm{A}_{\text{closed}} = \bm{A} + \bm{B}\bm{K}$，
其中$\bm{K}$为反馈增益矩阵。SAS对各模态的影响可以总结如下：

\begin{itemize}
    \item \textbf{短周期}：俯仰阻尼增益$k_q$增大俯仰阻尼，提高短周期阻尼比。
          姿态增益$k_\theta$增大等效静稳定性（类似增加$C_{m\alpha}$的绝对值），
          提高短周期频率。两者共同作用使短周期从可能欠阻尼（开环$\zeta_{sp} \approx 0.3$--$0.5$）
          变为良好阻尼（闭环$\zeta_{sp} \approx 0.7$--$1.0$）；
    \item \textbf{长周期}：SAS对长周期通常影响较小，因为长周期的动力学
         主要是能量交换而非力矩平衡。但过高的$k_\theta$可能通过迎角-姿态耦合
         间接影响长周期的阻尼；
    \item \textbf{荷兰滚}：偏航阻尼增益$k_r$直接改善荷兰滚阻尼。
          前摆的偏航力矩是荷兰滚阻尼的强有效执行器，因为偏航力矩是荷兰滚振荡中
          角加速度的直接驱动项；
    \item \textbf{滚转收敛}：滚转阻尼增益$k_p$和姿态增益$k_\phi$加速滚转收敛。
          差速反扭的滚转力矩对滚转速率的反馈是SAS滚转阻尼的主要来源
          （对于无副翼的纯推力矢量构型）；
    \item \textbf{螺旋模态}：SAS对螺旋模态的影响通过多个通道间接实现。
          偏航阻尼增益$k_r$对螺旋模态的影响取决于具体的飞行器参数——在某些情况下
          可能稳定螺旋模态，在其他情况下可能略微降低其稳定性。
          由于螺旋模态已在开环下基本可接受，SAS对它的影响通常不作为设计约束。
\end{itemize}

\begin{figure}[H]
\centering
\begin{tikzpicture}[>=Stealth, scale=1.0,
  annot/.style={font=\footnotesize\sffamily},
  sp/.style={fill=red!70!black, circle, inner sep=1.8pt, label={[font=\tiny\sffamily, red!70!black, align=center]}},
  ph/.style={fill=blue!70!black, circle, inner sep=1.8pt, label={[font=\tiny\sffamily, blue!70!black, align=center]}},
  dr/.style={fill=orange!80!black, circle, inner sep=1.8pt, label={[font=\tiny\sffamily, orange!80!black, align=center]}},
  rs/.style={fill=teal!70!black, circle, inner sep=1.8pt, label={[font=\tiny\sffamily, teal!70!black, align=center]}},
  spiral/.style={fill=purple!70!black, circle, inner sep=1.8pt, label={[font=\tiny\sffamily, purple!70!black, align=center]}},
  arrow/.style={->, thick, >={Stealth[length=2mm]}},
]
  % ===== 纵向模态 (s-plane, upper half) =====
  \draw[->] (-9,0) -- (1.5,0) node[annot, right] {$\Re$};
  \draw[->] (0,0) -- (0,4.2) node[annot, above] {$\Im$ (rad/s)};
  \node[annot] at (-7.5,4.5) {\textbf{纵向模态}};

  % 短周期 (开环)
  \node[sp] (spo) at (-1.8,3.5) {};
  \node[font=\tiny\sffamily, red!70!black, align=center] at (-1.8,4.0) {短周期（开环）\\$\omega_{sp}\approx2.5,\zeta\approx0.4$};
  % 短周期 (闭环 SAS)
  \node[sp] (spc) at (-4.2,2.2) {};
  \draw[arrow, red!50] (spo) -- (spc);
  \node[font=\tiny\sffamily, red!70!black, align=center] at (-4.5,2.6) {短周期（SAS闭环）\\阻尼增大、频率略增};

  % 长周期
  \node[ph] at (-0.15,0.55) {};
  \node[font=\tiny\sffamily, blue!70!black, align=center] at (0.8,0.7) {长周期\\$\omega_{ph}\approx0.5$\\$\zeta_{ph}\approx0.05$};
  \node[ph] at (-0.15,-0.55) {};

  % 标注轴
  \node[annot] at (-6.5,0.3) {$\leftarrow$ 更稳定};
  \draw[annot] (0.2,3.8) -- (0.6,3.8) node[right, font=\tiny\sffamily] {振荡频率};

  % ===== 横航向模态 (right side) =====
  \begin{scope}[shift={(8.5,0)}]
    \draw[->] (-6,0) -- (1.5,0) node[annot, right] {$\Re$};
    \draw[->] (0,0) -- (0,4.2) node[annot, above] {$\Im$ (rad/s)};
    \node[annot] at (-5,4.5) {\textbf{横航向模态}};

    % 滚转收敛 (开环)
    \node[rs] (rso) at (-2.8,0) {};
    \node[font=\tiny\sffamily, teal!70!black, align=center] at (-2.8,0.4) {滚转收敛\\$\tau\approx1$s};
    % 滚转收敛 (SAS)
    \node[rs] (rsc) at (-4.5,0) {};
    \draw[arrow, teal!50] (rso) -- (rsc);
    \node[font=\tiny\sffamily, teal!70!black] at (-5.2,0.4) {SAS: 更快收敛};

    % 荷兰滚 (开环)
    \node[dr] (dro) at (-0.6,3.2) {};
    \node[font=\tiny\sffamily, orange!80!black, align=center] at (0.2,3.5) {荷兰滚（开环）\\$\omega_{DR}\approx3,\zeta\approx0.15$};
    % 荷兰滚 (SAS)
    \node[dr] (drc) at (-3.0,2.5) {};
    \draw[arrow, orange!50] (dro) -- (drc);
    \node[font=\tiny\sffamily, orange!80!black] at (-3.8,2.8) {SAS: 阻尼增强};

    % 螺旋模态
    \node[spiral] at (0.08,0) {};
    \node[font=\tiny\sffamily, purple!70!black, align=center] at (0.8,-0.35) {螺旋模态\\$\lambda_S\approx\pm0.02$\\缓慢发散/收敛};
  \end{scope}
\end{tikzpicture}
\caption{SAS闭环对飞行模态特征根的影响（定性s平面示意图）。
上方：纵向模态——短周期（红色）在SAS作用下阻尼比显著提高（$0.4\to0.8$），
长周期（蓝色）受SAS影响很小。下方：横航向模态——滚转收敛（青色）加速，
荷兰滚（橙色）阻尼增强，螺旋模态（紫色）靠近虚轴且受SAS影响有限。
箭头定性指示SAS增益增大时极点的迁移方向。}
\label{fig:stability_poles}
\end{figure}

% ===========================================================================
\section{增稳控制律设计}
\label{sec:control}
% ===========================================================================

本节基于前文的动力学模型和控制效能分析，设计纵列双发矢量推力构型的增稳与控制律。
控制律采用级联结构，底层为SAS增稳内环，中层为姿态保持/角速度指令外环，
顶层为导航/轨迹规划（不在本文范围内）。

\subsection{级联控制架构}

控制架构分为三个明确的功能层次：

\begin{enumerate}[label=(\arabic*)]
    \item \textbf{操纵输入层}：人类操纵手（通过遥控器）或上层自主控制器
          （轨迹规划器/导航控制器）产生四通道操纵输入——
          油门$\omega_0$（纵向推力）、俯仰指令$\delta_t^{\text{pilot}}$、
          偏航指令$\delta_f^{\text{pilot}}$、滚转指令$\Delta\omega^{\text{pilot}}$；
    \item \textbf{SAS增强层}（Stability Augmentation System）：
          这是内环控制的核心，直接作用于角速率和姿态角反馈。
          SAS的使命是改善短周期阻尼、抑制大气湍流扰动、
          消除配平常值误差（通过积分项），并对飞行员操纵输入进行增稳修正；
    \item \textbf{执行层}：修正后的指令经物理限幅（摆角$\pm\delta_{\max}$，
          差速$\pm\Delta\omega_{\max}$）和电机动态模型后驱动实际执行器。
\end{enumerate}

\textbf{架构特征}：当前设计为\textbf{直接映射反馈增稳}，控制分配采用对角效能矩阵的
直接求逆（式\ref{eq:direct_mapping}），而非伪逆优化或二次规划在线求解。
这一简化依赖于第\ref{sec:allocation}节证明的通道天然解耦性质，
在正交摆轴几何约束下是完全合理的——可视为几何设计简化算法的经典案例。

\begin{figure}[H]
\centering
\begin{tikzpicture}[>=Stealth, node distance=8mm and 12mm,
  box/.style={draw=black!70, fill=black!8, rounded corners=3pt, minimum width=22mm, minimum height=9mm, align=center, font=\footnotesize\sffamily},
  sbox/.style={draw=black!50, fill=black!4, rounded corners=2pt, minimum width=16mm, minimum height=7mm, align=center, font=\footnotesize\sffamily},
  arrow/.style={->, thick, >={Stealth[length=2.5mm]}},
  annot/.style={font=\footnotesize\sffamily},
  dashbox/.style={draw=black!30, densely dashed, rounded corners=6pt, inner sep=6pt},
]
  % ===== 操纵输入（左侧） =====
  \node[box, fill=blue!8, draw=blue!40] (pilot) {操纵输入\\$\omega_0,\delta_t^{\text{pilot}},\delta_f^{\text{pilot}},\Delta\omega^{\text{pilot}}$};
  \node[annot, above=0pt of pilot] {操纵手/自主控制器};

  % ===== SAS层 =====
  \node[box, fill=red!8, draw=red!40, right=of pilot, xshift=4mm] (sas) {SAS 增稳层\\角速率阻尼\\+姿态P+I\\+角速度闭环};

  % ===== 限幅 =====
  \node[sbox, right=of sas, xshift=4mm] (limit) {指令限幅\\$\pm\delta_{\max},\pm\Delta\omega_{\max}$};

  % ===== 执行器 =====
  \node[box, fill=green!8, draw=green!50, right=of limit, xshift=4mm] (act) {执行器\\电机一阶\\滞后+摆座};

  % ===== 推力映射 =====
  \node[box, right=of act, xshift=4mm] (map) {推力矢量\\六维映射\\式(\ref{eq:six_component})};

  % ===== 动力学 =====
  \node[box, fill=orange!8, draw=orange!60, right=of map, xshift=4mm] (dyn) {六自由度\\刚体动力学\\NE+四元数};

  % ===== 传感器反馈（下方） =====
  \node[box, fill=black!5, draw=black!40, below=of sas, yshift=-6mm] (sensor) {传感器反馈\\$\bm{\omega}, \bm{q}, \bm{\theta}$};

  % ===== 扰动 =====
  \node[annot, above=of dyn, yshift=4mm] (dist) {气动扰动/湍流};

  % 前向连接
  \draw[arrow] (pilot) -- (sas);
  \draw[arrow] (sas) -- (limit);
  \draw[arrow] (limit) -- (act);
  \draw[arrow] (act) -- (map);
  \draw[arrow] (map) -- (dyn);
  \draw[arrow] (dist) -- (dyn);

  % 反馈连接
  \draw[arrow] (dyn.south) -- ++(0,-0.4) -| (sensor.east) node[pos=0.6, annot, right=2pt] {状态};
  \draw[arrow] (sensor.west) -| (sas.south) node[pos=0.35, annot, below=2pt] {反馈};

  % SAS 虚线框
  \begin{scope}[on background layer]
    \node[dashbox, fit=(sas)(sensor), label={[annot] above right:{\textbf{SAS 内环}}}] {};
  \end{scope}

  % 输出
  \draw[arrow] (dyn.east) -- ++(0.8,0) node[annot, right] {飞行状态};
\end{tikzpicture}
\caption{纵列双发矢量推力构型级联控制架构。三层结构：操纵输入层生成四通道指令，
SAS增强层通过角速率/姿态反馈修正指令，执行层经物理限幅后驱动电机和摆座。
控制分配采用对角效能矩阵的直接映射（式\ref{eq:direct_mapping}），无需在线优化求解。
内环反馈由IMU提供机体角速率和姿态四元数。}
\label{fig:control_architecture}
\end{figure}

\subsection{SAS基本增稳律}

三个通道的SAS反馈律以效率符号整定为原则。

\textbf{俯仰通道}（效率为负：$\partial M_y / \partial \delta_t = -b T_t < 0$）：
\begin{equation}
    \delta_t^{\text{cmd}} = \delta_t^{\text{pilot}}
    + k_q q + k_\theta \Delta\theta
    + k_i^\theta \int \Delta\theta \, dt
    \label{eq:sas_pitch}
\end{equation}
其中$k_q > 0$（效率为负$\to$正反馈使俯仰速率增大$\to$低头力矩增大$\to$抑制抬头速率）；
$k_\theta > 0$（俯仰角偏离越大$\to$摆角修正越大$\to$恢复力矩越大）；
积分项$k_i^\theta \int \Delta\theta \, dt$消除配平常值误差（如因推力线不完全过质心
或气动配平偏差导致的持续非零俯仰角）。

\textbf{偏航通道}（效率为正：$\partial M_z / \partial \delta_f = a T_f > 0$）：
\begin{equation}
    \delta_f^{\text{cmd}} = \delta_f^{\text{pilot}}
    - k_r r
    \label{eq:sas_yaw}
\end{equation}
效率为正意味着正摆角产生正偏航力矩（右偏航），因此偏航阻尼需要\textbf{负}反馈
（$-k_r r$）——当飞行器向右偏航（$r > 0$），控制律产生负摆角（左偏）
以产生逆方向的偏航力矩来抑制偏航速率。
偏航通道通常不包含姿态积分——航向保持是上层导航闭环的职责，
SAS仅需提供内环阻尼和角速率响应特性的改善。

\textbf{滚转通道}（效率为负：$\partial M_x / \partial \Delta\omega = -2k_Q \omega_0^2 < 0$）：
\begin{equation}
    \Delta\omega^{\text{cmd}} = \Delta\omega^{\text{pilot}}
    + k_p p + k_\phi \Delta\phi
    + k_i^\phi \int \Delta\phi \, dt
    \label{eq:sas_roll}
\end{equation}
效率为负$\to$正反馈，原理同俯仰通道。

\textbf{增益整定指南}（指导性原则，非标定数值）：
\begin{itemize}
    \item \textbf{角速率阻尼增益}（$k_q, k_r, k_p$）：从低值逐步增大，直至角速率阶跃响应
          的非周期阻尼（过阻尼）或目标阻尼比（0.7--1.0）。过高的角速率阻尼增益
          会引入高频噪声放大和执行器饱和风险；
    \item \textbf{姿态比例增益}（$k_\theta, k_\phi$）：通常设定为角速率阻尼增益的
          $\sim 1/4 \cdot \tau_{\text{姿态}}$量级，其中$\tau_{\text{姿态}}$为目标姿态
          响应时间常数；
    \item \textbf{积分增益}（$k_i^\theta, k_i^\phi$）：应足够低，使积分饱和
          恢复时间$\sim 5$--$10$秒。积分增益过高是诱发PIO（驾驶员诱发振荡）的常见原因。
\end{itemize}

\subsection{角速度闭环模式（Rate Command）}

角速度闭环模式（Rate Command / Rate CAC）是一种常见的飞行品质增强模式，
其中操纵杆（或滑块）的偏转直接映射为目标角速度，而非执行器位置指令。
松手后操纵杆弹簧回中、目标角速度归零、飞行器保持当前姿态（速率归零后的姿态）。
这种控制模式在遥控飞行器（如穿越机acro模式）和电传飞行器中广泛应用。

采用纯P控制器（比例控制）追踪目标角速率：
\begin{align}
    \delta_t^{\text{cmd}} &= k_q^{\text{rate}} (q - q_{\text{ref}}),
    &&\partial M_y / \partial \delta_t < 0 \Rightarrow (q - q_{\text{ref}})
    \text{为正时} \delta_t^{\text{cmd}} \text{取正} \\
    \delta_f^{\text{cmd}} &= k_r^{\text{rate}} (r_{\text{ref}} - r),
    &&\partial M_z / \partial \delta_f > 0 \Rightarrow (r - r_{\text{ref}})
    \text{为正时} \delta_f^{\text{cmd}} \text{取负} \\
    \Delta\omega^{\text{cmd}} &= k_p^{\text{rate}} (p - p_{\text{ref}}),
    &&\partial M_x / \partial \Delta\omega < 0 \Rightarrow (p - p_{\text{ref}})
    \text{为正时} \Delta\omega^{\text{cmd}} \text{取正}
    \label{eq:rate_command}
\end{align}

\textbf{性能分析}：纯P控制的角速度闭环存在稳态误差——
为维持非零角速率，需要非零的执行器偏转来克服气动阻尼力矩。
稳态误差可通过积分项消除，但积分项的引入会降低系统相位裕度，
在快速机动中可能导致超调和振荡。对于遥控飞行器（操纵手在闭环中），
适度的稳态误差是可以接受的——操纵手会自然地补偿并通过视觉反馈微调。

角速率闭环的跟踪带宽由P增益和电机响应共同决定。开环传递函数近似为
$G(s) = k^{\text{rate}} \cdot (\partial M / \partial \delta) \cdot
(1/(I s + |L_p|)) \cdot (1/(\tau_m s + 1))$，闭环带宽$\sim k^{\text{rate}}
\cdot |\partial M / \partial \delta| / I$（当电机响应足够快时）。

\subsection{饱和管理与积分器抗饱和}

所有执行器指令须经物理限幅，模拟真实舵机和电机的行程/速率限制：
\begin{equation}
    |\delta_t^{\text{cmd}}| \leq \delta_{\max}, \quad
    |\delta_f^{\text{cmd}}| \leq \delta_{\max}, \quad
    |\Delta\omega^{\text{cmd}}| \leq \Delta\omega_{\max}
    \label{eq:actuator_limits}
\end{equation}

积分器采用独立限幅和条件积分策略：
\begin{itemize}
    \item 积分器输出限幅$|\int \theta| \leq I_\theta^{\max}$防止姿态误差长时间持续
          导致的积分项无限增长（积分饱和，integrator windup）；
    \item 模式切换时积分器归零——当SAS模式从角速率闭环（mode 3）切换回全SAS（mode 1）时，
          旧模式下累积的积分值（基于不同误差信号）直接清零，防止在新模式刚启动时产生
          瞬态冲击（此问题在工程实践中被发现并修复）。
\end{itemize}

\subsection{差速分配与最大转速保护}

$\Delta\omega^{\text{cmd}}$由式(\ref{eq:diff_def})的平方根分配产生两电机目标转速。
此分配具有保持$\omega_f^2 + \omega_t^2 = 2\omega_0^2$常数、使一阶近似总推力不变的
优越性质。

然而在极限工况（高油门+大差速）下，一侧电机的目标转速可能超过最大允许转速$\omega_{\max}$。
在此情况下，必须施加电机级硬限幅：
\begin{equation}
    \omega_f^{\text{target}} \leftarrow \min(\omega_f^{\text{target}}, \omega_{\max}), \quad
    \omega_t^{\text{target}} \leftarrow \min(\omega_t^{\text{target}}, \omega_{\max})
    \label{eq:max_clamp}
\end{equation}

转速硬限幅的副作用是：\textbf{被钳位的电机转速低于目标值，导致平方转速和偏离$2\omega_0^2$，
总推力低于预期}。此外，由于只有一侧电机被钳位，差速指令的滚转力矩也不再精确等于
$-2k_Q\omega_0^2\Delta\omega^{\text{cmd}}$。SAS应在遥测中标记$w_{\max}$饱和状态，
提醒操纵手执行器已达物理极限，当前的控制指令无法被完全执行。

% ===========================================================================
\section{讨论：构型优势、局限与扩展路径}
\label{sec:discussion}
% ===========================================================================

\subsection{构型核心优势}

\begin{enumerate}[label=(\arabic*)]
    \item \textbf{通道天然解耦}：正交摆轴设计使控制效能矩阵退化为对角阵。
          这是构型设计层面的根本优势——解耦不是通过算法（如动态逆或鲁棒解耦）
          实现的，而是通过几何设计固化的。对角效能矩阵意味着每个控制通道仅需
          标定一个增益（$\partial M_i / \partial u_i$），无需在线矩阵求逆、
          无需伪逆优化、无需担心$\bm{B}(\omega_0)$随工作点退化的条件数问题；
    \item \textbf{推力矢量全速域有效}：推力矢量产生的控制力矩仅取决于推力
          大小和摆角，与空速无关。这颠覆了气动舵面的根本依赖——
          在悬停和极低速段（$V \approx 0$），气动舵面的控制力矩趋于零，
          而推力矢量的控制力矩完全保留。这一特性使本构型具备STOL甚至VTOL
          的潜力（前提是总推力$>$总重量，或配合适当的升力面）；
    \item \textbf{反扭主动对消与滚转控制一体化}：前后反向旋转螺旋桨使额定工况
          下反扭矩相互抵消——这不仅消除了不需要的寄生滚转力矩，还将滚转通道
          从“被动抑制”转变为“主动控制”。在这一点上，本构型的滚转控制
          比同轴直升机更简洁（同轴直升机需要复杂的上下旋翼总距差动），
          比四旋翼更高效（四旋翼通过改变四台电机的推力产生滚转力矩，
          是“间接控制”；本构型通过反扭矩差动直接产生滚转力矩，
          是“直接控制”）；
    \item \textbf{执行器数量最小化}：两个单轴摆座舵机+两台电机=四个执行器，
          对应于四个控制自由度（推力+三轴力矩），恰好达到完全可控所需的最少
          执行器数量。与四旋翼（4个电机，4个执行器）相当，但多出了固定翼的
          升力面；与六旋翼/八旋翼等过驱动多旋翼相比，执行器数量减半；
    \item \textbf{无尾翼/减尾翼潜力}：由于俯仰和偏航力矩均由推力矢量提供，
          传统的水平尾翼和垂直尾翼在理论上是冗余的。去尾翼意味着更少的气动阻力、
          更轻的结构重量、更小的雷达截面。实际工程中可保留小尾翼作为气动阻尼
          的补充和部分冗余。
\end{enumerate}

\subsection{已知局限与设计权衡}

任何构型设计都是在矛盾的需求之间寻找工程上可行的折中。以下列出本构型的主要局限，
它们不是缺陷，而是设计决策的内在权衡——了解这些权衡对于在正确的应用场景中
使用这一构型至关重要。

\begin{enumerate}[label=(\arabic*)]
    \item \textbf{低油门滚转控制退化}：差速滚转效能$\propto \omega_0^2$，
          在低油门时退化显著。对于需要精细滚转控制的飞行阶段（如进近着陆、
          编队飞行），这一退化可能构成安全风险。缓解措施包括：
          (a) 限制飞行包线下界的油门开度；(b) 保留小型气动副翼作为低速段的
          滚转控制补充；(c) 在着陆构型中使用襟翼增大升力以允许更低进近速度
          同时保持油门安全阈值；
    \item \textbf{$w_{\max}$饱和导致推力偏离}：大差速+高油门组合可能导致
          一侧电机超速。转速硬限幅虽然保护了电机，但破坏了推力不变假设。
          这会在机动最剧烈时产生额外的总推力扰动；
    \item \textbf{俯仰-偏航耦合不可完全消除}：反扭矩的摆角投影耦合
          在同时大偏转工况下可达主控力矩的$20\%$--$60\%$。
          通过前馈补偿可显著降低，但在极限机动中（如协调转弯的剧烈俯仰+偏航组合），
          耦合仍可能产生可感知的交叉轴扰动；
    \item \textbf{双推进系统全部失效即失去全部控制}：与常规固定翼
          （可通过气动舵面滑翔迫降）和多旋翼（即使一台电机失效，剩余电机
          可通过重构继续可控飞行\cite{ducard2009modeling}）不同，
          纯推力矢量构型在两台电机同时失效后丧失所有姿态控制力矩。
          发动机空中重启能力和供电系统冗余是安全底线要求；
    \item \textbf{欧拉角万向锁}：3-2-1欧拉角在$\theta = \pm 90^\circ$处奇异。
          这限制了控制律在欧拉角域中的适用性——任何以欧拉角为反馈变量的控制律
          在$\theta \to \pm 90^\circ$时都会遇到奇异性。
          采用四元数直接反馈（如误差四元数$\times$角速率反馈）可以
          消除这一问题\cite{markley2014fundamentals}，但增加了控制律复杂度；
    \item \textbf{推力方向余弦损失}：摆角偏转使推力有效前向分量降低$\cos\delta$倍。
          $\delta = 25^\circ$满偏时前向推力损失约$10\%$（$\cos 25^\circ \approx 0.906$），
          这一损失在持续机动（如持续大摆角盘旋）时累计可观。
          在性能计算和任务规划中应计入这一损失。
\end{enumerate}

\subsection{与现有构型的综合比较}

\begin{table}[H]
\centering
\caption{纵列双发矢量推力构型与现有典型构型的定性比较}
\begin{tabular}{p{2.8cm}p{3cm}p{3cm}p{3cm}p{3cm}}
\toprule
\textbf{特性} & \textbf{纵列双发矢量推力} & \textbf{常规固定翼} & \textbf{四旋翼} & \textbf{倾转旋翼} \\
\midrule
低速控制权限 & 高（推力矢量） & 低$\to$无（舵面） & 高（差速） & 高（推力矢量） \\
巡航效率 & 高（固定翼升力面） & 最高（优化翼型） & 低（全桨盘升力） & 中（倾转机构重量） \\
控制分配复杂度 & 极低（对角B） & 低（经典混控） & 低（混控矩阵） & 高（过驱动B$^+$） \\
机械复杂度 & 低（2舵机+2电机） & 低（3-4舵机） & 低（4电机） & 高（倾转机构） \\
VTOL/STOL能力 & 可能（T/W $>$ 1） & 需跑道 & 是 & 是 \\
执行器失效冗余 & 无（双发全部失效=失控） & 气动舵面备份 & 部分（6+旋翼） & 部分（倾转冗余） \\
万向锁风险 & 存在（欧拉角反馈） & 低（常规包线） & 低（常规包线） & 存在（倾转过渡） \\
\bottomrule
\end{tabular}
\end{table}

\begin{figure}[H]
\centering
\begin{tikzpicture}[>=Stealth, scale=0.82,
  config/.style={draw=black!60, fill=black!6, rounded corners=3pt, minimum width=32mm, minimum height=22mm, align=center, font=\footnotesize\sffamily},
  arrow/.style={->, thick, >={Stealth[length=2mm]}},
]
  \node[config, draw=blue!50, fill=blue!5] (a) at (0,0) {
    \textbf{常规固定翼}\\[2pt]
    {\small\color{black!30}固定推力方向}\\[2pt]
    舵面控制 $\cdot$ 需跑道\\[1pt]
    \textcolor{green!50!black}{巡航效率最高}\\
    \textcolor{red!60!black}{低速控制弱}\\
    \textcolor{red!60!black}{无法悬停}
  };

  \node[config, draw=green!50, fill=green!5] (b) at (5.2,0) {
    \textbf{四旋翼}\\[2pt]
    {\small\color{black!30}差速控制}\\[2pt]
    差速控制 $\cdot$ 悬停\\[1pt]
    \textcolor{green!50!black}{低速操纵优异}\\
    \textcolor{red!60!black}{巡航效率低}\\
    \textcolor{yellow!50!black}{续航受限}
  };

  \node[config, draw=orange!60, fill=orange!5] (c) at (10.4,0) {
    \textbf{倾转旋翼}\\[2pt]
    {\small\color{black!30}推力方向可变}\\[2pt]
    倾转过渡 $\cdot$ VTOL\\[1pt]
    \textcolor{green!50!black}{全速域覆盖}\\
    \textcolor{red!60!black}{机构复杂沉重}\\
    \textcolor{red!60!black}{过渡控制难}
  };

  \node[config, draw=red!60, fill=red!3] (d) at (2.6,-6.0) {
    \textbf{本文构型：纵列双发矢量推力}\\[2pt]
    {\small\color{red!50}正交摆轴 $\cdot$ 天然解耦}\\[2pt]
    矢推+差速反扭 $\cdot$ 天然解耦\\[1pt]
    \textcolor{green!50!black}{全速域可控 $\cdot$ 执行器最少}\\
    \textcolor{yellow!50!black}{低油门滚转退化}\\
    \textcolor{green!50!black}{通道天然解耦 $\cdot$ 无在线优化}
  };

  % 箭头指示本构型的定位
  \draw[arrow, blue!60, line width=1.2pt] (a.south) -- ++(0,-1.2) node[midway, right, font=\footnotesize, align=center] {无矢推\\需跑道};
  \draw[arrow, green!60, line width=1.2pt] (b.south) -- ++(0,-1.2) node[midway, right, font=\footnotesize, align=center] {差速控制\\低效巡航};
  \draw[arrow, orange!60, line width=1.2pt] (c.south) -- ++(0,-1.2) node[midway, right, font=\footnotesize, align=center] {矢推VTOL\\机构复杂};

  % 收敛到本构型
  \draw[arrow, red!80, line width=1.4pt] (2.6,-3.8) -- (2.6,-5.0) node[midway, right, font=\footnotesize\bfseries, align=center] {融合优势\\规避劣势};
\end{tikzpicture}
\caption{纵列双发矢量推力构型在现有构型谱系中的定位。该构型融合了固定翼的巡航效率
（升力面提供主要升力）、多旋翼的差速滚转机制（反扭矩差动）和倾转旋翼的全速域
推力矢量控制能力，同时通过正交摆轴的几何设计规避了倾转机构复杂和控制分配在线优化
的代价。图中定性地展示了各种构型在机械复杂度、巡航效率和低速控制三个维度上的权衡。}
\label{fig:config_comparison}
\end{figure}

\subsection{技术成熟度与验证路径}
\label{sec:verification}

本构型目前处于概念设计与仿真验证阶段，距飞行试验还有相当的技术距离。
建议的技术成熟度推进路径为：

\textbf{Tier 1 — 参数标定}：在推力台架上测量$k_T(\omega)$、$k_Q(\omega)$、
$\tau_m(\omega)$函数关系，标定电机-桨-ESC系统模型。同步进行CFD计算，
获取气动导数数据库（$C_L(\alpha,\beta)$、$C_D(\alpha,\beta)$、
$C_m(\alpha,\beta,\delta_t)$等）。

\textbf{Tier 2 — 硬件在环（HIL）}：将标定后的参数模型和SAS控制律部署到
实时仿真平台（如Pixhawk HIL或dSPACE），连接真实的飞控硬件（IMU、舵机、电调），
验证控制律的实时性和执行器接口的正确性。

\textbf{Tier 3 — 系留飞行}：在有安全绳约束的条件下进行有限自由度的飞行试验。
系留绳索提供垂直方向约束和断电紧急回收，允许在真实气动环境中验证悬停能力
和基本姿态控制。

\textbf{Tier 4 — 自由飞行}：逐步扩展飞行包线，从悬停、低速平飞到巡航，
最后验证大机动和边界工况。

每一步应设置明确的通过/不通过判据，并记录与仿真预测的偏差以回溯更新模型。

\subsection{未来扩展方向}

\begin{enumerate}[label=(\arabic*)]
    \item \textbf{控制分配升级}：从直接映射升级为加权伪逆+零空间优化，
          在保持主通道对角解耦的同时，利用剩余自由度实现耦合项前馈补偿
          和执行器负载均衡。特别是当引入气动舵面后，执行空间由3扩至6+，
          控制分配的必要性凸显；
    \item \textbf{过渡走廊探索}：研究从悬停（可能需要$\pm 90^\circ$摆座）
          到平飞巡航的连续过渡策略。过渡段的气动-推力耦合是最具挑战的控制问题——
          倾转旋翼的过渡控制已有大量研究\cite{kang2009control}，可借鉴其增益调度
          和模型预测控制（MPC）思路。对于纵列双发构型，由于推力线始终在机身平面上，
          过渡中的推力-升力耦合比倾转旋翼更平缓，可能是一个意外优势；
    \item \textbf{全气动-推力融合}：在中高速巡航段引入气动舵面，
          形成推力矢量（全速域）+舵面（高效巡航+冗余）的混合控制架构。
          混合控制的关键在于设计合理的控制分配策略，在舵面权限充足时优先使用舵面
          （减小推力损失和摆座磨损），在舵面不足时无缝过渡到推力矢量；
    \item \textbf{非线性控制}：在大机动工况（摆角$>25^\circ$，小角度近似失效）
          下，线性SAS的性能退化。考虑引入模型预测控制（MPC）、
          滑模控制或反馈线性化等非线性控制方法，在发挥矢量推力机动潜力的同时
          保证稳定性；
    \item \textbf{故障容错}：设计执行器失效后的控制重构策略。对于电机单发失效，
          可设计降级控制律（利用剩余的单发电机的推力矢量和反扭矩维持部分控制），
          将飞行器引导至安全迫降。对于舵机失效，可设计锁定摆座+备用电机的应急策略。
\end{enumerate}

% ===========================================================================
\section{结论}
\label{sec:conclusion}
% ===========================================================================

本文系统阐述了纵列双发正交单轴矢量推力飞行器构型的理论基础，按从推进系统物理到
整机控制律的逻辑链展开。主要工作包括：

\begin{enumerate}[label=(\arabic*)]
    \item \textbf{推进系统建模}：从桨盘动量理论和叶素理论出发，建立了推力和反扭矩的
          集总参数模型（$T = k_T\omega^2$，$\tau = k_Q\omega^2 + J_p\dot{\omega}$），
          阐述了量纲分析的相似律（$k_T \propto \rho D^4$，$k_Q \propto \rho D^5$），
          并分析了电机一阶动态和旋转部件陀螺效应的物理机制；
    \item \textbf{推力矢量映射}：通过摆座运动学分析，导出了完整的六维力/力矩映射
          （式\ref{eq:six_component}），分离了主控通道与耦合通道，
          揭示了推力线过质心和正交摆轴是确保解耦性的必要几何条件；
    \item \textbf{六自由度动力学}：基于牛顿-欧拉方程和四元数姿态运动学，
          建立了整机刚体动力学模型，给出了子步调度的数值积分策略；
    \item \textbf{控制分配分析}：证明了在小角度近似下，控制效能矩阵退化为严格对角阵
          （式\ref{eq:effectiveness_matrix}），正交摆轴设计使控制分配从欠定优化
          简化为三通道直接映射。定量评估了耦合项的量级（约为推力力臂项的$D/a$倍）；
    \item \textbf{稳定性模态分析}：通过小扰动线性化，分析了纵向两模态（短周期、
          长周期）和横航向三模态（滚转收敛、荷兰滚、螺旋）的物理本质，
          以及SAS闭环对模态特性的影响；
    \item \textbf{控制律设计}：给出了基于效率符号整定的三通道SAS增稳律、
          角速度闭环P控制律，以及饱和管理和抗积分饱和策略。
\end{enumerate}

该构型的核心创新在于：\textbf{通过几何设计（正交摆轴、推力线过质心、反向旋转对消）
实现了控制通道的天然解耦}，使控制分配问题从通用的欠定优化退化为简单的直接映射。
这一设计哲学——用几何换取算法简洁性——在小型电推进飞行器的工程约束下具有实际价值：
它使飞控开发更简单（无需在线优化求解器），使可靠性分析更直接（通道独立、故障隔离清晰），
使调参更直观（每个增益仅影响一个通道的响应特性）。

本文的分析不绑定具体飞行器物理参数，重在揭示构型本身的动力学本质与控制特性。
文中导出的解析公式（式\ref{eq:six_component}的六维映射、式\ref{eq:effectiveness_matrix}
的效能矩阵、式\ref{eq:coupling_ratio}的耦合比）可作为参数优化、控制律设计
和子系统选型的理论框架。对于具体型号的工程实现，还需在参数标定（台架试验）、
气动数据获取（CFD或风洞）、实时仿真验证（HIL）和飞行试验等多个维度持续迭代。

% ===========================================================================
%  参考文献
% ===========================================================================
\begin{thebibliography}{99}

\bibitem{stevens2015aircraft}
Stevens B L, Lewis F L, Johnson E N.
\textit{Aircraft Control and Simulation: Dynamics, Controls Design, and Autonomous Systems}.
3rd ed. Hoboken, NJ: Wiley, 2015.

\bibitem{leishman2006helicopter}
Leishman J G.
\textit{Principles of Helicopter Aerodynamics}.
2nd ed. Cambridge: Cambridge University Press, 2006.

\bibitem{mccormick1995aerodynamics}
McCormick B W.
\textit{Aerodynamics, Aeronautics, and Flight Mechanics}.
2nd ed. New York: Wiley, 1995.

\bibitem{ducard2011fault}
Ducard G J J.
\textit{Fault-tolerant Flight Control and Guidance Systems: Practical Methods for Small Unmanned Aerial Vehicles}.
London: Springer, 2011.

\bibitem{johansen2013control}
Johansen T A, Fossen T I.
Control allocation --- A survey.
\textit{Automatica}, 2013, 49(5): 1087--1103.

\bibitem{markley2014fundamentals}
Markley F L, Crassidis J L.
\textit{Fundamentals of Spacecraft Attitude Determination and Control}.
New York: Springer, 2014.

\bibitem{kuipers1999quaternions}
Kuipers J B.
\textit{Quaternions and Rotation Sequences: A Primer with Applications to Orbits, Aerospace and Virtual Reality}.
Princeton, NJ: Princeton University Press, 1999.

\bibitem{etkin1982dynamics}
Etkin B.
\textit{Dynamics of Atmospheric Flight}.
New York: Wiley, 1982.

\bibitem{stengel2004flight}
Stengel R F.
\textit{Flight Dynamics}.
Princeton, NJ: Princeton University Press, 2004.

\bibitem{beard2012small}
Beard R W, McLain T W.
\textit{Small Unmanned Aircraft: Theory and Practice}.
Princeton, NJ: Princeton University Press, 2012.

\bibitem{mahony2012multirotor}
Mahony R, Kumar V, Corke P.
Multirotor aerial vehicles: Modeling, estimation, and control of quadrotor.
\textit{IEEE Robotics \& Automation Magazine}, 2012, 19(3): 20--32.

\bibitem{bresciani2008quadrotor}
Bresciani T.
Modelling, identification and control of a quadrotor helicopter.
MSc Thesis, Lund University, 2008.

\bibitem{ding2018modeling}
Ding C, Wang C, Lu L, Ouyang B.
Modeling and control of fully actuated vector thrust unmanned aerial vehicles.
In: \textit{Proceedings of 2018 International Symposium on Flexible Automation (ISFA)},
Kanazawa, Japan, 2018: Paper ISFA2018-L034.

\bibitem{seguigasco2014novel}
Segui-Gasco P, Al-Rihani Y, Shin H S, Savvaris A.
A novel actuation concept for a multi rotor UAV.
\textit{Journal of Intelligent \& Robotic Systems}, 2014, 74(1): 173--191.

\bibitem{zhu2025experimental}
Zhu H, Du Y, Nie H, Xin Z, Geng X.
Experimental investigation of aerodynamic interaction in non-parallel tandem dual-rotor
systems for tiltrotor UAV.
\textit{Drones}, 2025, 9(5): 374.

\bibitem{oppenheimer2010control}
Oppenheimer M W, Doman D B, Bolender M A.
Control allocation for over-actuated systems.
In: \textit{Proceedings of the 14th Mediterranean Conference on Control and Automation},
2006: 1--6.

\bibitem{kang2009control}
Kang Y, Hedrick J K.
Linear tracking for a fixed-wing UAV using nonlinear model predictive control.
\textit{IEEE Transactions on Control Systems Technology}, 2009, 17(5): 1202--1210.

\bibitem{ducard2009modeling}
Ducard G, Geering H P.
Efficient nonlinear actuator fault detection and isolation system for unmanned aerial vehicles.
\textit{Journal of Guidance, Control, and Dynamics}, 2008, 31(1): 225--237.

\bibitem{kim2004nonlinear}
Kim H J, Shim D H.
A flight control system for aerial robots: algorithms and experiments.
\textit{Control Engineering Practice}, 2003, 11(12): 1389--1400.

\end{thebibliography}

\end{document}
